\subsection*{Proof structure}
The finite compiler and balanced cubic routing of Part~I turn the sparse
system into a local network with controlled residual transport;
Lemma~\ref{c:lem:logical} states the interface used here. The physical construction supplies a
local Bloch-band minimum, scalar defects and interaction estimates that
realize this network. A finite-box inverse bound, including the
infinite-dimensional gradient space, controls the remaining physical modes.
Finite material, source and readout approximations complete the reduction.

\subsection*{Why a local band minimum is sufficient}
The periodic medium constructed below has a simple minimum on a small
Bloch chart. The selected band is the top of the sixteen-fold corner
cluster of Section~\ref{c:sec:edge}, and its minimum is local: the
quadratic regime extends only over a chart of radius $q_0=\varrho2^{-40}$
in Bloch momentum. A defect of radius $L$ must be localized in momentum
inside this chart, which forces the defect radius to satisfy
$L\ge\Ls$ in~\eqref{c:eq:masterconstants}. No assertion is made that its energy lies in a complete
three-dimensional gap. A cutoff resolvent on that chart is an
\emph{auxiliary proof operator}, not an implemented material or a physical
filter. The auxiliary response differs from a genuine Maxwell response
by a quantified $L^2$ residual. The full finite-box inverse bound controls
that difference. This distinction allows all three spatial directions to
be used without assuming an invariant physical sector.

\section{A full-operator inverse bound on any finite PEC box}
\label{c:sec:detuning}

This section is independent of the particular dielectric platform.

\begin{lemma}[The finite spectral count]
\label{c:lem:count}
Let $\Omega=\prod_{j=1}^3(0,N_j)$ with integer $N_j\ge1$, and let
$1\le\varepsilon_0\le25$ be any bounded measurable real scalar function.
For $A_0=\Cc-\varepsilon_0$, there are at most
\begin{equation}
 N_\Omega=2\prod_{j=1}^3(1+2N_j)\le54|\Omega|             \label{c:eq:count}
\end{equation}
spectral values in $[-1/4,1/4]$, counting multiplicity. These are isolated
finite-multiplicity eigenvalues; there is no essential spectrum in that
interval.
\end{lemma}
\begin{proof}
Use the component families
\[
 \begin{array}{ll}
 E_1:&\cos(i_1\pi r_1/N_1)\sin(i_2\pi r_2/N_2)\sin(i_3\pi r_3/N_3),\\
 E_2:&\sin(i_1\pi r_1/N_1)\cos(i_2\pi r_2/N_2)\sin(i_3\pi r_3/N_3),\\
 E_3:&\sin(i_1\pi r_1/N_1)\sin(i_2\pi r_2/N_2)\cos(i_3\pi r_3/N_3).
 \end{array}
\]
The cosine index may be zero; the two sine indices are positive. After
normalization these are a complete orthonormal component basis of $L^2$.
Grouping by $(i_1,i_2,i_3)$ diagonalizes the curl form into the matrices
$|p|^2I-pp^T$, with fixed harmless component signs. With three
positive indices there is one zero direction and two eigenvalues
$|p|^2$; with one zero index there is one positive direction.

Here is a domain justification. Integration by parts against a compactly
supported smooth vector field and Parseval show convergence of its block
expansion in the curl-form norm. Every finite expansion has zero
tangential trace. Its zero extension has no surface term in its curl;
inward dilation and mollification approximate it in the form norm by
compactly supported smooth fields, i.e., compactly supported smooth fields
are dense in the fields with zero tangential trace~\citep[Chapter~3]{monk2003}. Taking closures therefore gives
exactly $H_0(\curl;\Omega)$, not a smaller smooth subspace.

Let $P_{\rm g}$ and $P_{\rm t}$ be the orthogonal projections onto $\ker\Cc$ and its
orthogonal complement. The restriction $\Cc_{\rm t}$ to the complement has compact resolvent, with positive eigenvalues
\[
 \pi^2\sum_{j=1}^3(\nu_j/N_j)^2,
\]
with at least two positive indices and multiplicity at most two per
triple. Write the bounded material blocks as
$\varepsilon_{\rm gg}=P_{\rm g}\varepsilon_0P_{\rm g}\ge I$, $\varepsilon_{\rm tt}=P_{\rm t}\varepsilon_0P_{\rm t}\le25I$, and
$\varepsilon_{\rm tg}=P_{\rm t}\varepsilon_0P_{\rm g}$. The compression of $A_0$ to $\ker\Cc$ is
$-\varepsilon_{\rm gg}$, with spectrum in $[-25,-1]$, so the window $[-1/4,1/4]$ keeps
distance $3/4$ from it. For $|\lambda|\le1/4$, the gradient block
$-(\varepsilon_{\rm gg}+\lambda)$ is invertible. Its Schur operator is
\begin{equation}
 S_{\rm t}(\lambda)=\Cc_{\rm t}-\varepsilon_{\rm tt}-\lambda+\varepsilon_{\rm tg}(\varepsilon_{\rm gg}+\lambda)^{-1}\varepsilon_{\rm tg}^*.
                                                               \label{c:eq:transverseschur}
\end{equation}
It has domain $D(\Cc_{\rm t})$, compact resolvent, and
$S_{\rm t}(\lambda)\ge\Cc_{\rm t}-26I$. Moreover, for $\lambda_2>\lambda_1$,
\[
 S_{\rm t}(\lambda_2)\le S_{\rm t}(\lambda_1)-(\lambda_2-\lambda_1)I.
\]
Its ordered eigenvalues are continuous by the bounded-perturbation
min--max inequality and strictly decrease in this sense. Each crosses
zero at most once. Only as many of them as there are eigenvalues of $\Cc_{\rm t}$ at most $26$ can do so, because the others
are positive throughout the interval. Schur factorization shows that
these crossings are exactly the spectrum of $A_0$ in the interval, with
the same nullity; at every other point it gives a bounded inverse into
$D(\Cc)$. This also excludes essential spectrum there.

Finally $\sqrt{26}/\pi<2$, so each index of a free eigenvalue at most
26 is less than $2N_j$. Counting triples with multiplicity two gives
\eqref{c:eq:count}. The infinite gradient space has not been counted as a
finite collection of modes; it was eliminated as a uniformly invertible
operator block.
\end{proof}

\begin{lemma}[Uniform scalar detuning]
\label{c:lem:detuning}
For $0<\tau\le1/64$ dyadic, let $J_{\rm grid}$ be the least power of two at least
$64mN_\Omega$. Choose $j$ uniformly in $\{0,\ldots,J_{\rm grid}-1\}$ and put
\begin{equation}
 \sigma=\frac{(j+1/2)\tau}{J_{\rm grid}},\qquad
 L_b=\frac{64mN_\Omega}{\tau}.                         \label{c:eq:detuning}
\end{equation}
Provided the material interval is preserved, $A=A_0-\sigma I$ satisfies
\eqref{c:eq:maininverse}, apart from its final comparison
$702L_b\le m^{c_A}$, except on an event of probability at most
$3/(64m)$. That comparison is proved in Section~\ref{c:sec:cost}.
\end{lemma}
\begin{proof}
An interval of radius $L_b^{-1}$ around one real eigenvalue contains at
most $2J_{\rm grid}/(L_b\tau)+1$ grid points. Spectrum outside $[-1/4,1/4]$ is
farther than $L_b^{-1}$ from the detuning interval. Lemma~\ref{c:lem:count}
and a union bound give
\[
 \Pr\{\dist(\sigma,\spec A_0)<L_b^{-1}\}
 \le\frac{2N_\Omega}{L_b\tau}+\frac{N_\Omega}{J_{\rm grid}}
 \le\frac3{64m}.
\]
Singular draws are included; equality at the endpoint is harmless.
Self-adjointness gives $\norm{A^{-1}}_{2\to2}\le L_b$ on the
complementary event. If $E=A^{-1}f$, then
$\norm{\Cc E}\le\norm f+25\norm E$, whence
$\norm{A^{-1}}_{2\to X}\le27L_b$ and
$\norm A_{X\to2}<26$. Their product is at most $702L_b$.
\end{proof}

The form representation in~\eqref{c:eq:domain} uses the first representation
theorem~\citep[Chapter~VI, Section~2, Theorem~2.1]{kato1976}. Its hypotheses
hold because $V$ is dense in $L^2$, complete in the form norm, and the
form is nonnegative and symmetric. Bounded real multiplication preserves
self-adjointness on its operator domain. The displayed Fourier argument
also proves the required constant-coefficient decomposition directly.

\section{An explicit scalar periodic medium with a simple local minimum}
\label{c:sec:edge}

Set
\begin{equation}
 \varrho=2^{-100},\qquad
 V_0(r)=\cos(\pi r_1)+2\cos(\pi r_2)+4\cos(\pi r_3),
 \qquad \varepsilon_{\rm per}=7+\varrho V_0.           \label{c:eq:background}
\end{equation}
The period is two in each coordinate. At the antiperiodic corner
$k_*=(\pi/2,\pi/2,\pi/2)$, use local Bloch momentum $q$ and Fourier
aliases $p_\nu(q)=k_*+q+\pi\nu$, $\nu\in\Z^3$. The free fiber is the direct
sum of
\begin{equation}
 |p_\nu|^2I-p_\nu p_\nu^T-7I.                                \label{c:eq:freefiber}
\end{equation}
For complex momentum the squares and transpose in this formula are
bilinear, not Hermitian. Real momenta give the self-adjoint fibers.
The potential $V_0$ connects only neighboring aliases, with weights
$1/2,1,2$.

For complex $q$, the free fiber is the maximal multiplication operator
on $\ell^2(\Z^3;\C^3)$ with domain
\[
 \left\{u:\sum_\nu\big\|[p_\nu(q)^2I-p_\nu(q)p_\nu(q)^T-7I]u_\nu\big\|^2
                    <\infty\right\}.
\]
It is closed, since convergence of a sequence and of its image gives
the defining identity in each finite alias block. The material
convolution is bounded, so its addition uses this same domain.
For real $q$ this is the self-adjoint Bloch fiber obtained by the
unitary Fourier decomposition of the physical curl form. The Bloch
transform groups ordinary Fourier frequencies modulo $\pi\Z^3$;
Parseval proves this direct-integral decomposition, initially on
smooth compactly supported fields and then by graph closure.

We do not assume that these unbounded operator domains are independent
of momentum. All complex inverses used below are constructed from the
explicit free inverse
\begin{equation}
 (p^2I-pp^T-\mathcal EI)^{-1}
   =\frac{I-pp^T/\mathcal E}{p^2-\mathcal E},                           \label{c:eq:freeinverse}
\end{equation}
and a bounded Neumann series. Multiplying the block formulas verifies
both inverse identities; the image belongs to the maximal domain
because the product with the free operator is the identity plus a
bounded operator. Uniform bounds on the alias blocks and on the
Neumann series prove holomorphy in the bounded-operator norm.
This constructs the complex inverse and finite Schur operators
without a common-domain perturbation assumption.

\subsection{The corner cluster}
At $q=0$, the eight wavevectors $p_\varsigma=\tfrac\pi2\varsigma$, $\varsigma\in\{\pm1\}^3$,
have a sixteen-dimensional transverse eigenspace at
$E_0=3\pi^2/4-7$. Let $P_\varsigma=I-\varsigma\varsigma^T/3$ and let $P_{\rm cl}$ be the projection onto this cluster. The first-order perturbation on
this space is $K_{\rm cl}=-P_{\rm cl}V_0P_{\rm cl}$. Reflection parity $\varpi\in\{\pm1\}^3$ reduces
it to the plane $\one^\perp\subset\R^3$. With
$b_\varpi=(\varpi_1,2\varpi_2,4\varpi_3)$ and $P_1=I-\one\one^T/3$, the block is
\begin{equation}
 K_{\rm cl}(\varpi)=P_1\diag(b_\varpi)P_1-\frac{\varpi_1+2\varpi_2+4\varpi_3}{2}P_1.
                                                               \label{c:eq:parityblock}
\end{equation}
To verify this, an amplitude $\phi\in\one^\perp$ at $\varsigma=(1,1,1)$ has
amplitude $\varpi^{\iota(\varsigma)}\diag(\varsigma)\phi/\sqrt8$ at $\varsigma$, where
$\iota_j(\varsigma)=(1-\varsigma_j)/2$. Toggling $\varsigma_j$ multiplies this amplitude by the
polar reflection $I-2e_je_j^T$ and parity $\varpi_j$. Applying the coupling
$-\varpi_j2^{j-2}$ gives~\eqref{c:eq:parityblock} after projection. Distinct parity embeddings are orthogonal, because
the sum of a nontrivial sign character on $\{\pm1\}^3$ is zero.
Each projected parity space has dimension two, so the eight spaces
exhaust the sixteen-dimensional transverse cluster.

The two roots of this block are
\begin{equation}
 -\frac{\varpi_1+2\varpi_2+4\varpi_3}{6}
 \ \,\pm\,\frac{\sqrt{21-2\varpi_1\varpi_2-4\varpi_1\varpi_3-8\varpi_2\varpi_3}}3.
                                                               \label{c:eq:clusterroots}
\end{equation}
The unique largest root is
\begin{equation}
 \lambda_*=(5+2\sqrt{19})/6,
 \quad \varpi=(1,-1,-1),\qquad
 \lambda_*-\lambda_{\rm next}>1/20.                  \label{c:eq:clustergap}
\end{equation}
The next root is $1/2+\sqrt3$. Substitution in the eight displayed
quadratics proves the ordering; the accompanying exact program checks
all sixteen intervals and also rebuilds the $24\times24$ corner
projection and its parity intertwiners.

Let $u_0$ be the normalized top cluster vector. Its eight wavevector
blocks have equal norm by reflection symmetry. Projecting the free
velocity onto the transverse cluster gives
\[
 \Upsilon(t)=\pi\diag_{\varsigma}(\varsigma\cdot t)I_{\varsigma^\perp}.
\]
Consequently
\begin{equation}
 \ip{u_0}{\Upsilon(t)u_0}=0,\qquad
 \norm{\Upsilon(t)u_0}^2=\pi^2|t|^2.                    \label{c:eq:velocity}
\end{equation}
If $\lambda_K(t)$ is the top eigenvalue of $K_{\rm cl}+\Upsilon(t)$ near zero, ordinary
finite-dimensional Schur elimination gives
\begin{equation}
 t^TH_0t=\ip{\Upsilon(t)u_0}{(\lambda_*-K_{\rm cl})^{-1}_{u_0^\perp}\Upsilon(t)u_0},
 \qquad H_0=\tfrac12\lambda_K''(0).
                                                               \label{c:eq:A0}
\end{equation}
Since $\norm{K_{\rm cl}}\le\norm{V_0}_\infty\le7$, the spectrum of $K_{\rm cl}$ lies in
$[-7,7]$, so $(\lambda_*-K_{\rm cl})^{-1}\ge1/14$ on $u_0^\perp$. The gap
\eqref{c:eq:clustergap} gives $(\lambda_*-K_{\rm cl})^{-1}\le20$ there.
With~\eqref{c:eq:velocity} this gives
\begin{equation}
 \frac{\pi^2}{14}I\le H_0\le20\pi^2I.           \label{c:eq:A0bounds}
\end{equation}
The exact cluster roots lie in $[-2.29,2.29]$, so the lower constant is
far from sharp. Reflections make $H_0$
diagonal. Its exact entries are
\begin{equation}
 H_0=\pi^2\diag
 \left(\frac{47}{31}+\frac{352\sqrt{19}}{589},\ 
 2+\frac{11\sqrt{19}}{38},\ 
 \frac18+\frac{11\sqrt{19}}{152}\right).              \label{c:eq:explicitA0}
\end{equation}
For a direct verification, take
$\phi^{(0)}_i=(\lambda^{\rm d}_i-\lambda_*)^{-1}$ with $\lambda^{\rm d}=(7/2,1/2,-3/2)$, normalize $\phi^{(0)}$,
and evaluate~\eqref{c:eq:A0} in the parity block with the corresponding
sign flipped. This uses only arithmetic in $\Q(\sqrt{19})$.

\subsection{The continuum remainder, including every Fourier alias}
\begin{lemma}[Local continuum band]
\label{c:lem:edge}
The actual continuum operator $A_{\rm per}=\curl\curl-\varepsilon_{\rm per}$
has a simple analytic Bloch eigenvalue $\lambda(q)$ for
$|q_j|<R_q:=\varrho2^{-12}$. It is real for real $q$ and even in each
coordinate. With $E_c=\lambda(0)$ and $H=\tfrac12\lambda''(0)$,
\begin{align}
 |E_c-E_0-\varrho\lambda_*|&\le2^{16}\varrho^2,\\
 \norm{\varrho H-H_0}&\le2^{40}\varrho,\\
 (2\varrho)^{-1}I&\le H\le(256/\varrho)I.              \label{c:eq:edgebounds}
\end{align}
For $|q|\le R_q/2$,
\begin{equation}
 |\lambda(q)-E_c-q^THq|\le M_4|q|^4,
 \qquad M_4=4\varrho R_q^{-4}.                       \label{c:eq:quartic}
\end{equation}
On the smaller chart $|q_j|\le2q_0$, $q_0=\varrho2^{-40}$, the selected
band is separated from every other fiber spectral value by at least
$\varrho/100$. The separated inverse and eigenprojection are bounded
holomorphic operators. The selected cell eigenfunction has Fourier
$\ell^1$ norm at most 32 in a complex neighborhood of this chart.
\end{lemma}
\begin{proof}
Put $q=\varrho t$ and $\mathcal E=E_0+\varrho e$. Retain all three components
of the eight central aliases. Every other free alias has transverse
energy at least $11\pi^2/4-7$, while its longitudinal denominator at
energy $\mathcal E$ is $-3\pi^2/4-\varrho e$. Formula~\eqref{c:eq:freeinverse}
shows that its inverse norm is less than one for
$|t_j|\le2^{-12}$ and $|e|\le8$. Indeed the transverse denominator is
at least $2\pi^2-O(\varrho|p|)-8\varrho$; its ratio to the
quadratic numerator is uniformly bounded, also as $|p|\to\infty$.
Using $3<\pi<4$ gives the stated conservative bound. The eight central
longitudinal directions have the same bound, since their diagonal
free correction is $q^2-(s\cdot q)^2/3$.

The potential has norm at most seven. Its complementary inverse is the
Neumann series with ratio at most $7\varrho<1/4$, and has norm at most
two. The cluster--complement coupling, divided by $\varrho$, has norm
at most eight. Eliminating the complement therefore gives the exact
sixteen-dimensional analytic Schur matrix
\begin{equation}
 K_{\rm cl}+\Upsilon(t)-eI+\varrho\Sigma_2(t,e),\qquad \norm{\Sigma_2(t,e)}\le256.
                                                               \label{c:eq:feshbach}
\end{equation}
The bound is $2\cdot8^2$ for the self-energy plus less than one for
the direct quadratic term. This derivation includes the infinitely many
longitudinal aliases; none was discarded by a transverse approximation.

On $|e-\lambda_*|=1/100$,
$\norm{\Upsilon(t)}\le3\pi2^{-12}<1/400$, so the resolvent of $K_{\rm cl}+\Upsilon(t)$ is
less than 140. Its contour projection $\Pi(t)$ has rank one and norm at
most two. Hence
$(K_{\rm cl}+\Upsilon(t)-e)^{-1}=\Pi(t)/(\lambda_K(t)-e)+\widetilde R(e)$, where the reduced
resolvent $\widetilde R(e)=(K_{\rm cl}+\Upsilon(t)-e)^{-1}(I-\Pi(t))$ is analytic inside
the circle. On the circle its norm is at most $140\cdot3$, so by the
maximum principle it is at most $420$ inside. This bounds the resolvent
near its root by $2/|e-\lambda_K(t)|+2048$. On the circle
$|e-\lambda_K(t)|=2^{16}\varrho$, multiplication by
$\varrho\Sigma_2$ has norm less than $1/16$. Homotopy of the finite
Schur determinant, or its contour projection, thus leaves exactly one
root and gives
$|e_\varrho(t)-\lambda_K(t)|\le2^{16}\varrho$.
The root is analytic because its contour projection has rank one.
The complementary free inverse proves that this root is an actual
continuum fiber eigenvalue. It is simple. The homotopy counts zeros of
the Schur determinant with multiplicity, so the root is a simple zero,
and the inverse Schur matrix (adjugate over determinant) has a simple
pole with rank-one residue. The Feshbach formula writes the fiber
resolvent as this inverse between two analytic lifting operators plus an
analytic term; the Riesz projection of the full fiber therefore has rank
one.

Cauchy's estimate on the $t$-polydisc gives the Hessian error
$2^{40}\varrho$. The finite inequalities in~\eqref{c:eq:A0bounds}
then give~\eqref{c:eq:edgebounds}. Coordinate reflections of the scalar
material, with the polar-vector action on fields, send $q_j$ to $-q_j$.
Simplicity gives the asserted evenness. The analytic bound
$|\lambda(q)-E_0|\le3\varrho$ and the one-variable Cauchy estimate
on each complex line give~\eqref{c:eq:quartic}: after removing the
constant and quadratic terms, the even tail is at most
$3\varrho(|q|/R_q)^4/[1-(|q|/R_q)^2]$.

The cluster gap persists on the smaller chart by the same Schur
estimates. The other free spectrum stays a fixed distance away.
For the cell eigenvector, the complementary Neumann series starts in
finitely many neighboring aliases. Each potential application changes
an alias by one, and its $\ell^1$ norm is at most seven. The free
inverse has norm at most one on each noncentral $3\times3$ block.
The geometric series, together with the eight central blocks, bounds
the eigenvector's Fourier $\ell^1$ norm by 32. Cauchy bounds its
derivatives. This proves all assertions without a finite-element
spectral convergence assumption.
\end{proof}

For later normalization choose the analytic cell eigenfunction
$\Phi_q$ with mean squared norm one for real $q$, real $\Phi_0$,
and dual pairing $\langle\Phi_{-q},\Phi_q\rangle_{\rm bil}=1$.
One construction is to apply the analytic rank-one projection to
$\Phi_0$ and divide by the analytic square root of its nonzero pairing
with $\Phi_0$. The transpose identity $A_{\rm per}(q)^T=A_{\rm per}(-q)$ in the
antiperiodic cell gauge gives $\Pi(-q)=\Pi(q)^T$ for the analytic
eigenprojection, and hence this normalization. The pairing stays away
from zero: on $|t_j|\le2^{-39}$ the projection moves by $O(2^{-33})$
from its value at $q=0$, where the pairing is one. At purely imaginary
$q$, both the operator and the normalized eigenfunction are real.
This fact will fix the sign of the useful coupling numerator.

A real zeroth-order field, used only as a starting vector for calibration,
is
\begin{equation}
 \Phi^{(0)}(r)=\sqrt8\begin{pmatrix}
 \phi^{(0)}_1\sin(\pi r_1/2)\sin(\pi r_2/2)\sin(\pi r_3/2)\\
 -\phi^{(0)}_2\cos(\pi r_1/2)\cos(\pi r_2/2)\sin(\pi r_3/2)\\
 -\phi^{(0)}_3\cos(\pi r_1/2)\sin(\pi r_2/2)\cos(\pi r_3/2)
 \end{pmatrix},                                      \label{c:eq:initialphi}
\end{equation}
where $\phi^{(0)}$ is the normalized vector following~\eqref{c:eq:explicitA0}.
It is divergence-free, has mean squared norm one, and the selected
actual field can be signed so that
$\norm{\Phi_0-\Phi^{(0)}}_{\ell^1({\rm Fourier})}\le2^{16}\varrho$.
This follows by the same isolated-cluster Schur equation and geometric
complement series.

\subsection{A fixed integer placement matrix}
Use
\begin{equation}
 B=\diag(7180532256056349,\ 6388184103869743,\ 2347449745487106).
                                                               \label{c:eq:B}
\end{equation}
The exact identities~\eqref{c:eq:explicitA0}, rational enclosures for $\pi$
and $\sqrt{19}$, and~\eqref{c:eq:edgebounds} imply
\begin{equation}
 .99\,|y|\le|H^{-1/2}By|\le1.01\,|y|.               \label{c:eq:Bmetric}
\end{equation}
The verifier checks the stronger squared inequalities, allowing the full
$2^{-60}=2^{40}\varrho$ Hessian error. This is an integer placement
lattice, not a change to the Euclidean metric or to material tensors.

\section{Smooth cutoffs and the auxiliary resolvent}
\label{c:sec:cutoff}

Define the fixed bump and its square root by
\begin{equation}
 w(y)=\begin{cases}\exp[-1/(1-|y|^2)],&|y|<1,\\0,&|y|\ge1,\end{cases}
 \qquad q_w(y)=\sqrt{w(y)}.                           \label{c:eq:well}
\end{equation}
For a one-dimensional cutoff, normalize
$\psi_0(s)=\exp[-1/(1-4s^2)]\one_{|s|<1/2}$ to have integral one and put
$\psi=\one_{[-3/2,3/2]}*\psi_0$. Its normalization integral is between
$1/8$ and one. Thus $0\le\psi\le1$, $\psi=1$ on $[-1,1]$, and
$\psi=0$ outside $[-2,2]$. In the corner Bloch coordinates set
\begin{equation}
 \chi(q)=\prod_{j=1}^3\psi(q_j/q_0).                \label{c:eq:chi}
\end{equation}
It is a scalar multiplier on the Bloch torus, extended by zero outside
this chart. In physical Fourier variables it is periodic modulo the
reciprocal lattice. In particular, it commutes with every translation
and every constant-coefficient derivative.

\begin{lemma}[Quantitative Gevrey estimates]
\label{c:lem:gevrey}
The zero extensions in~\eqref{c:eq:well} are smooth. For every directional
derivative of order $K$,
\begin{equation}
 \norm{\partial^K w}_\infty+\norm{\partial^K q_w}_\infty
 \le2\cdot1024^K(K!)^2.                              \label{c:eq:gevreyder}
\end{equation}
They satisfy the Fourier tail bound
\begin{equation}
 \int_{|p|\ge R}(|(\mathcal Fw)(p)|+|(\mathcal Fq_w)(p)|)\dd p
 \le2^{100}\exp(-\sqrt R/512).                       \label{c:eq:gevreytail}
\end{equation}
The cutoff $\chi$ has derivative bounds
$2^{45}(2^{16}/q_0)^K(K!)^2$. All its values, and the values of $w,q_w$,
can be enclosed to $N_{\rm b}$ bits in time polynomial in $N_{\rm b}$.
\end{lemma}
\begin{proof}
For a real point inside the unit ball, let $s_y=1-|y|^2$.
In the complex polydisc of radius $s_y/100$, the real part of
$(1-z^2)^{-1}$ is at least $1/(2s_y)$; for the square root use
$1/(4s_y)$ in the exponential. Cauchy's estimate gives
$K!(100/s_y)^K e^{-1/(4s_y)}$. Maximizing in $s_y$ and using
$K^K e^{-K}\le K!$ bounds coordinate derivatives by
$400^K(K!)^2$. The directional expansion gives~\eqref{c:eq:gevreyder}
with the displayed larger constant. These estimates also show that every
derivative tends to zero at the boundary, justifying the smooth zero
extension.

Integrate by parts $K$ times in the direction of $p$ and use the
support volume less than eight. Choosing an integer
$K$ within one of $\sqrt{|p|}/64$ gives
$|(\mathcal Fw)(p)|+|(\mathcal Fq_w)(p)|\le
2^{20}e^{-\sqrt{|p|}/256}$. For small $|p|$ use the $L^1$
bound. Radial integration and one further weakening of the exponent
prove~\eqref{c:eq:gevreytail}. Rescaling the one-dimensional bump,
dividing by its normalization at least $1/8$, and taking the three
products gives the claimed cutoff bounds.

For evaluation, if the positive denominator in the bump is less than
$1/[4(N_{\rm b}+10)]$, the value is below $2^{-N_{\rm b}-2}$ and can be enclosed directly.
Otherwise the exponential argument is $O(N_{\rm b})$. Its Taylor series with
range reduction and a geometric factorial-tail bound uses polynomially
many rational operations and bits. A periodic trapezoidal rule applied
to a smooth compactly supported bump has error bounded by the sum of its
Fourier coefficients at the aliased indices;~\eqref{c:eq:gevreytail}
therefore gives $O(N_{\rm b}^2)$ nodes per dimension for fixed intervals.
This encloses the normalization and the convolution defining $\psi$.
Equivalently, one may integrate a polynomial approximation with the same
Gevrey tail. Outward rational rounding gives certified intervals.
\end{proof}

For clarity, the following fixed constants are used throughout:
\begin{equation}
 \Cs=2^{2^{40}},\qquad \cs=2^{-1024},\qquad
 c_u=2^{-1024},\qquad
 \Ls=8^{\lceil64\cdot2^{40}/3\rceil},\qquad
 \ds=\Cs^8.                                         \label{c:eq:masterconstants}
\end{equation}
Their size is only a convenience for a single majorant ledger. The proofs
below give bounds far smaller than these numbers. Section~\ref{c:app:ledger}
records the derivative and arithmetic majorants that justify their use;
none is a supremum over inaccessible data.

Let $L\ge\Ls$ be a power of eight and $h=L^{-1}$. Compute a dyadic $z_h$
with
\begin{equation}
 |z_h-(E_c-h^2)|\le h^4.                              \label{c:eq:z}
\end{equation}
Section~\ref{c:sec:calibration} gives a certified terminating algorithm.
For $0\le\sigma\le h^4/64$, define on the whole space
\begin{equation}
 R_\sigma=\chi(q)\big(A_{\rm per}(q)-z_h-\sigma\big)^{-1}
 \quad\text{on the chart, and zero elsewhere}.         \label{c:eq:R}
\end{equation}
This inverse exists on the chart by Lemma~\ref{c:lem:edge}. It is a real
self-adjoint bounded operator, with
\begin{equation}
 \norm{R_\sigma}\le4L^2,\quad
 \norm{R_\sigma-R_0}\le16\sigma L^4,\quad
 (A_{\rm per}-z_h-\sigma)R_\sigma=\chi.               \label{c:eq:Ridentities}
\end{equation}
The last equality holds as an operator into $L^2$, and
$R_\sigma L^2\subset D(\curl\curl)$ on $\R^3$.
The norm estimates follow from the selected denominator, at least
$h^2/2$, and the fixed separation of all other bands. The denominator
bound uses $\lambda(q)\ge E_c$ on the chart: there
$M_4|q|^2\le M_4(4q_0)^2<10^{-6}/\varrho$, far below the smallest
eigenvalue $(2\varrho)^{-1}$ of $H$, so~\eqref{c:eq:quartic} gives
$\lambda(q)-E_c\ge q^THq/2$. The final factor
$\chi$ in~\eqref{c:eq:Ridentities} can not be replaced by $I$; this is
the main analytic issue of the proof.

\section{A broad scalar defect with an isolated auxiliary mode}
\label{c:sec:defect}

Let $\mathbb B_L=\{r:|r|<L\}$, let $Q_h$ multiply by $q_w(hx)$ and restrict
to $\mathbb B_L$, and define on $L^2(\mathbb B_L;\C^3)$
\begin{equation}
 S_h=Q_hR_0Q_h^*.                                    \label{c:eq:Sh}
\end{equation}
The adjoint includes the same multiplier and extension by zero. This is
an infinite-dimensional self-adjoint operator; it is not declared compact.

\begin{lemma}[An effective scalar operator with a quantitative gap]
\label{c:lem:perron}
Let $H$ satisfy~\eqref{c:eq:edgebounds}. The scalar compact operator
\begin{equation}
 T_H=\sqrt w\,(-\nabla\cdot H\nabla+1)^{-1}\sqrt w
 \quad\text{on }L^2(\mathbb B_1)                              \label{c:eq:TA}
\end{equation}
has a simple largest eigenvalue $\mu_H$ and a real positive normalized
eigenvector $u_H$. It satisfies
\begin{equation}
 2^{-127}\le\mu_H\le2^{-92},\qquad
 \dist(\mu_H,\spec T_H\setminus\{\mu_H\})\ge2^{-190}.
                                                               \label{c:eq:perrongap}
\end{equation}
In addition,
\begin{equation}
 u_H(y)\ge2^{-29}\sqrt{w(y)},\qquad
 \int\sqrt w\,u_H\ge2^{-32}.                         \label{c:eq:perronoverlap}
\end{equation}
\end{lemma}
\begin{proof}
Fourier transformation gives the Yukawa kernel
\[
 Y_H(y)=\frac{e^{-|H^{-1/2}y|}}
 {4\pi\sqrt{\det H}\,|H^{-1/2}y|}.
\]
It can also be obtained by changing variables in the radial fundamental
solution of $-\Delta+1$, whose normalization follows by integrating its
normal derivative over a small sphere. On $\mathbb B_1\times\mathbb B_1$ the bounds
on $H$ imply
\begin{equation}
 Y_H(r-y)\ge Y_-:=\varrho2^{-24},\quad
 Y_H(r-y)\le\frac{4\varrho}{|r-y|},\quad
 \sup_r\norm{Y_H(r-\cdot)}_{L^2(\mathbb B_1)}\le Y_+:=32\varrho.
                                                               \label{c:eq:yukawabounds}
\end{equation}
For example, $|H^{-1/2}(r-y)|<1$, so the exponential is at least
$1/4$; inserting the extreme eigenvalues of $H$ proves the first
inequality. Squaring the second and integrating over a radius-two ball
gives the third. The Schur integral bound gives $\norm{T_H}\le256\varrho$.
Also $\int w\ge1/8$: on the half-radius ball, $w\ge e^{-4/3}>1/4$,
and its volume is $\pi/6$.

The kernel is square integrable on $\mathbb B_1\times\mathbb B_1$, so $T_H$ is
compact. It is nonnegative as a quadratic form and has a strictly
positive kernel in the interior. Its largest Rayleigh quotient is
attained; replacing an eigenvector by its absolute value does not
reduce it, and the integral equation makes the resulting eigenvector
strictly positive. Testing on $\sqrt w$ gives
$\mu_H\ge Y_-\int w\ge\varrho2^{-27}=2^{-127}$.
The integral equation and Cauchy--Schwarz give
\[
 u_H(r)\le(Y_+/\mu_H)\sqrt{w(r)},\qquad
 \int\sqrt w\,u_H\ge\mu_H/Y_+,
 \qquad u_H(r)\ge(Y_-/Y_+)\sqrt{w(r)}.
\]
These imply~\eqref{c:eq:perronoverlap}.

For the spectral gap define the reversible Markov kernel
\[
 \mathcal K(r,y)=\frac{\sqrt{w(r)}Y_H(r-y)\sqrt{w(y)}u_H(y)}
 {\mu_Hu_H(r)}.
\]
It has invariant measure $u_H^2\dd r$. The preceding upper bound on
$u_H$ implies
$\mathcal K(r,y)\ge\varsigma_0u_H(y)^2$, where
$\varsigma_0=Y_-\mu_H/Y_+^2\ge2^{-61}$.
Write $\mathcal K=\varsigma_0\Pi_{\rm c}+(1-\varsigma_0)\mathcal K_1$, where $\Pi_{\rm c}$ is the orthogonal projection onto constants and $\mathcal K_1$ is another
Markov operator preserving that measure. Jensen's inequality makes
$\mathcal K_1$ a contraction on $L^2(u_H^2\dd r)$. On the mean-zero subspace,
$\norm{\mathcal K}\le1-\varsigma_0$. The unitary map $v\mapsto v/u_H$
identifies $T_H/\mu_H$ with $\mathcal K$. This proves simplicity and a gap at
least $\mu_H\varsigma_0\ge\varrho2^{-88}>2^{-190}$.
No uncomputed scalar eigenvalue is a premise of this argument.
\end{proof}

\begin{lemma}[Quantitative effective-mass comparison]
\label{c:lem:effectivemass}
For every $L\ge\Ls$, the largest eigenvalue $\mu_h$ of $S_h$ is simple,
real, and isolated. It and a real normalized eigenvector $u_h$ satisfy
\begin{align}
 2^{-128}L^2&\le\mu_h\le2^{-90}L^2,
 &\dist(\mu_h,\spec S_h\setminus\{\mu_h\})&\ge2^{-192}L^2,
                                                               \label{c:eq:mugap}\\
 \norm{u_h-h^{3/2}\Phi_0(r)u_H(hx)}_{L^2(\mathbb B_L)}&\le2^{921} h\le\Cs h.
                                                               \label{c:eq:modeapprox}
\end{align}
The sign can be chosen so that the normalized starting vector
$h^{3/2}q_w(hx)\Phi^{(0)}(r)$ has overlap at least $2^{-48}$ with $u_h$.
Moreover,
\begin{equation}
 \norm{u_h}_{H^1(\R^3)}\le2^{512}L^2\le\Cs L^2,                   \label{c:eq:modeH1}
\end{equation}
where zero extension is understood.
\end{lemma}
\begin{proof}
Split $R_0$ into its selected-band part and the remaining part. The latter
has norm at most $2^{112}$ by Lemma~\ref{c:lem:edge}; after multiplication
by $h^2$ it is a bounded perturbation of norm at most $2^{112}h^2$.

Write the selected normalized Bloch eigenvector in alias coordinates as
$\phi(q)$. Define an isometry from the scaled momentum chart by
\[
 (\mathcal U_h v)(r)=h^{3/2}(2\pi)^{-3/2}
 \int e^{ihp\cdot r}\Phi_{hp}(r)v(p)\dd p,
\]
and put
\[
 s_h(p)=\frac{h\sqrt{\chi(hp)}}{\sqrt{\lambda(hp)-z_h}},
 \qquad s_H(p)=(1+p^THp)^{-1/2}.
\]
Extend $s_h$ by zero off the chart. The selected part of $h^2S_h$ is
$B_hB_h^*$, where $B_h=Q_h\mathcal U_hs_h$. Its nonzero spectrum is
that of $B_h^*B_h$.

The latter operator has a concrete Fourier kernel. Up to the unitary
Fourier factor $(2\pi)^{-3/2}$, its middle factor is
\begin{equation}
 \sum_{\nu,\nu'}\phi_\nu(hp)^*\phi_{\nu'}(hp')
 (\mathcal Fw)\big(p-p'+\pi(\nu-\nu')/h\big).                \label{c:eq:umklapp}
\end{equation}
For $\nu=\nu'$ this is the kernel of multiplication by $w$ conjugated by
Fourier transform, with the vector multipliers $\phi(hp)$ and $\phi(hp')$.
For $\nu\ne\nu'$, the argument of $\mathcal Fw$ is at least
$1/h$ in norm: both $hp,hp'$ lie in the very small fixed chart.
The two Fourier $\ell^1$ bounds $32$ on $\phi$, the pointwise bound
$|(\mathcal Fw)(p)|\le2^{20}e^{-\sqrt{|p|}/256}$ from the proof of
Lemma~\ref{c:lem:gevrey}, and the volume, less than $h^{-3}$, of the
scaled chart in the row and column Schur integrals make the norm of
this off-alias part at most
\begin{equation}
 2^{100}h^{-3}e^{-1/(512\sqrt h)}\le2^{512}h.           \label{c:eq:umklapperror}
\end{equation}
The last inequality follows by maximizing $s^8e^{-s/512}$ for $s\ge0$;
the stated constant is deliberately larger than that maximum.

The dispersion estimate gives
\begin{equation}
 \sup_p|s_h(p)-s_H(p)|\le2^{128}h,
 \qquad
 \sup_p s_H(p)\norm{\phi(hp)-\phi(0)}\le2^{128}h.            \label{c:eq:multipliererror}
\end{equation}
Here is an explicit check of the first assertion. On the chart the
normalized denominator differs from $1+p^THp$ by at most
$h^2+M_4h^2|p|^4$ and is at least one half of it. Applying the derivative
of $s^{-1/2}$ bounds the difference by
$C h^2(1+M_4|p|^4)/(1+2^{99}|p|^2)^{3/2}$.
For $|p|\le4q_0/h$ its supremum is less than $2^{100}h$.
Where the cutoff is not one, $s_H\le h/(q_0\sqrt{2^{99}})$,
which fits the displayed bound. The second assertion follows from the
Cauchy derivative bound on $\phi$, using
$|p|s_H(p)\le2^{-99/2}$. These calculations use only the explicit
powers in Lemma~\ref{c:lem:edge}.

Multiplication by $w$ has norm at most one. The diagonal-alias
comparison has norm at most $2^{133}h$, so adding the off-alias
bound in~\eqref{c:eq:umklapperror} gives
\begin{equation}
 \norm{B_h^*B_h-
 s_H\mathcal Fw\mathcal F^{-1}s_H}\le2^{513}h\le2^{520}h.
                                                               \label{c:eq:effectivecompare}
\end{equation}
The nonzero spectrum of the comparison operator is exactly that of
$T_H$. The selected part is compact: its momentum chart is bounded,
its multiplier is bounded for fixed $h$, and the spatial support of
$Q_h$ is bounded, so its integral kernel is Hilbert--Schmidt.
The separated part need not be compact, but its norm after scaling is
at most $2^{112}h^2$. The min--max inequalities above that bounded
essential spectrum, or the elementary isolated-eigenvalue Schur
argument, now transfer~\eqref{c:eq:perrongap} to~\eqref{c:eq:mugap}.
The choice $L\ge\Ls$ makes both errors smaller than $2^{-200}$, so the
top eigenvalue of $h^2S_h$ lies in $[\mu_H-2^{-199},\mu_H+2^{-199}]$
and the rest of its spectrum stays at distance at least
$2^{-190}-2^{-198}$ from it. Multiplying by $L^2=h^{-2}$
gives~\eqref{c:eq:mugap}.

For completeness, eigenvectors can be compared without identifying two
different Hilbert spaces. Let $y_H$ be the normalized top eigenvector
of $s_H\mathcal Fw\mathcal F^{-1}s_H$. Projection of the perturbed
eigenvector equation onto $y_H^\perp$ bounds its component there by
twice the error in~\eqref{c:eq:effectivecompare} divided by the gap.
Reconstruct the spatial eigenvector as $B_hy_h/\sqrt{\mu_hh^2}$.
The operator $B_h$ differs by $O(h)$, with the explicit bounds just
used, from
\[
 v\longmapsto h^{3/2}\Phi_0(r)\sqrt{w(hx)}
       [\mathcal F^{-1}(s_Hv)](hx).
\]
To check this operator comparison, use the disjoint Bloch aliases on
$|p|\le q_0/h$,~\eqref{c:eq:multipliererror}, and the bound
$\sup_{|p|>q_0/h}s_H(p)\le2^{100}h$ on the remaining input.
The comparison operator applied to $y_H/\sqrt{\mu_H}$ is exactly
$h^{3/2}\Phi_0(r)u_H(hx)$. Finally the bounded separated-band
perturbation is treated by the same gap projection. The scalar-well
argument above gives the stronger gap $2^{-188}$, so the conservative
Gram error $2^{520}h$ gives an eigenvector error at most $2^{710}h$.
On the eigenvalue interval bounded below by $2^{-128}$, the inverse
square root and its Lipschitz constant are both bounded by $2^{200}$.
Combining these estimates with the reconstruction error gives
$2^{913}h$ for the selected spatial mode. Adding the separated-band
perturbation gives $2^{920}h$; taking its real part and normalizing
increases the error by at most a factor of two. This proves
\eqref{c:eq:modeapprox}, with all constants below the majorant of
Section~\ref{c:app:ledger}.

The starting-vector overlap is positive quantitatively. By
\eqref{c:eq:perronoverlap}, its leading integral is at least
$2^{-29}\int w(y)|\Phi_0(y/h)|^2\dd y$, apart from the
$\Phi^{(0)}-\Phi_0$ error. The mean of $|\Phi_0|^2$ is one.
Its nonconstant Fourier coefficients are absolutely summable, and
\eqref{c:eq:gevreytail} shows that the integral is at least $1/16$.
The field error $2^{16}\varrho$, the mode error $\Cs h$, and
normalization leave the conservative overlap $2^{-48}$.

For~\eqref{c:eq:modeH1}, use translations rather than an assumed elliptic
regularity theorem for all polarizations. The cutoff $\chi$ commutes
with translations. The local inverse without its factor $\chi$ has
norm at most $4L^2$. Its resolvent identity under translation of
$\varepsilon_{\rm per}$ gives
\[
 \norm{\mathcal T_aR_0\mathcal T_a^{-1}-R_0}\le2^{20}L^4|a|.
\]
Together with the derivative bound on $q_w(hx)$ this gives
$\norm{\mathcal T_aS_h\mathcal T_a^{-1}-S_h}\le2^{32}L^4|a|$. Projecting $\mathcal T_au_h$ onto
$u_h^\perp$, using the gap $2^{-192}L^2$, and choosing the real sign
by positive overlap yields
$\norm{\mathcal T_au_h-u_h}\le2^{256}L^2|a|$ for small $a$.
The Fourier difference-quotient characterization then gives
$u_h\in H^1$ with~\eqref{c:eq:modeH1}. Explicitly, Fatou applied to
$|e^{iap_j}-1|^2/a^2$ bounds $\int p_j^2|(\mathcal Fu_h)(p)|^2\dd p$.
This also justifies the zero extension across the ball boundary.
\end{proof}

Define the compactly supported source mode
\begin{equation}
 a_h=Q_h^*u_h=q_w(hx)u_h(r).                          \label{c:eq:source-mode}
\end{equation}
It satisfies the exact identity
\begin{equation}
 a_h=\mu_h^{-1}w(hx)R_0a_h.                           \label{c:eq:source-mode-equation}
\end{equation}
This identity, not a claim of an exact isolated eigenstate of the full
infinite Maxwell operator, will control leakage.

\begin{lemma}[Nonzero analytic source numerator]
\label{c:lem:numerator}
For real $\beta$ with $|\beta|\le h$, define the two analytic band
source factors by
\begin{align*}
 F_h^+(q)&=(2\pi)^{-3/2}\int e^{-iq\cdot r}
             \Phi_{-q}(r)\cdot a_h(r)\dd r,\\
 F_h^-(q)&=(2\pi)^{-3/2}\int e^{iq\cdot r}
             \Phi_q(r)\cdot a_h(r)\dd r.
\end{align*}
At $q=i\beta$ both are real and
\begin{equation}
 h^{3/2}F_h^\pm(i\beta)\ge2^{-48}.                   \label{c:eq:numeratorlower}
\end{equation}
The scaled product
$\mathcal N_h(k)=h^3F_h^-(hk)F_h^+(hk)$ has uniformly bounded
analytic derivatives of every fixed order on the neighborhoods used
below. Up to order 32 the explicit bounds in
Section~\ref{c:app:ledger} apply, independently of $h$.
\end{lemma}
\begin{proof}
At imaginary momentum the cell functions are real by the normalization
following Lemma~\ref{c:lem:edge}. Insert~\eqref{c:eq:modeapprox}, put
$y=hx$, and use $e^{\pm\beta\cdot r}\ge e^{-1}$ on the support.
The leading term is bounded below by
\[
 (2\pi)^{-3/2}e^{-1}2^{-29}
 \int w(y)|\Phi_0(y/h)|^2\dd y.
\]
The last integral is at least $1/16$, as in the preceding proof.
The change from $\Phi_0$ to $\Phi_{\pm i\beta}$ is at most a fixed
explicit constant times $h$, by its Fourier $\ell^1$ Cauchy bound.
The $L^2$ mode error contributes at most $\Cs h$ times a fixed
support-volume and cell-field bound. Since $L\ge\Ls$, these errors
are smaller than half the displayed positive lower bound, which proves
\eqref{c:eq:numeratorlower}.

Differentiation in $k$ introduces factors $hx$ of norm at most one,
and derivatives of $\Phi_{hk}$ introduce $h$ times its fixed Cauchy
constants. Cauchy--Schwarz bounds the rescaled $L^1$ norm of $a_h$ by a
fixed number. Thus all fixed-order derivatives have the stated uniform
bounds, even for real $k$ as large as the scaled chart. No differentiability
of $u_h$ beyond its compact support is needed for this assertion.
\end{proof}

The scalar tuning interval is already explicit. At $v_* =1/\mu_h$,
$v_*^{-1}I-S_h$ has a simple zero mode and complementary separation
$2^{-192}L^2$. Every target diagonal in
$[-2^{-196}L^2,2^{-196}L^2]$ is realized by
$v=(\mu_h+\Theta_{ii})^{-1}$. These positive values are at most
$2^{130}L^{-2}<1$. This is the auxiliary defect resonance used by the
construction; its error as a physical Maxwell mode is established in
Section~\ref{c:sec:leakage}.

\section{Three-dimensional interactions with a strict rate margin}
\label{c:sec:green}

Translations below are ordinary Euclidean translations. For
$R_\nu=2LdB\nu$, $\nu\in\Z^3\setminus\{0\}$, define
\[
 G_h(\nu)=Q_hR_0\mathcal T_{R_\nu}Q_h^*,\qquad
 h_j(d)=\ip{u_h}{G_h(e_j)u_h}.
\]
Since $L$ is even, the corner phase $e^{ik_*\cdot R_\nu}$ equals one.

\begin{lemma}[Full blocks and useful scalar couplings]
\label{c:lem:interactions}
For $L\ge\Ls$ and every integer $d\ge1$,
\begin{equation}
 \norm{G_h(\nu)}\le\Cs L^5
       \big(3^{-d|\nu|_1}+e^{-\cs L}\big).              \label{c:eq:blockdecay}
\end{equation}
If $d\ge\ds$, $d^{15}\ge\Cs^4L$, and
\begin{equation}
 \Cs L^5e^{-\cs L}\le(c_u/4)L^2 8^{-d},              \label{c:eq:couplingtail}
\end{equation}
then, in all three coordinate directions,
\begin{equation}
 h_j(d)\ge c_uL^2 8^{-d}.                            \label{c:eq:couplinglower}
\end{equation}
The same upper bounds hold with $R_\sigma$ in place of $R_0$, after
adding the norm error $16\sigma L^4$ for the whole compressed union.
In particular the useful and unwanted exponential bases obey
$(1/3)^2<1/8$.
\end{lemma}
This margin separates couplings with $|\nu|_1\ge2$ from the useful ones.
Two occupied sites at $|\nu|_1=1$ would couple at the useful scale; the
parity layout of Section~\ref{c:sec:network} places only intended
terminal--mediator incidences at that distance. Condition
\eqref{c:eq:couplingtail} is guaranteed by the parameter tests
\eqref{c:eq:Ltailtest} and~\eqref{c:eq:tail-test}.

\subsection{Contour movement for a smooth, nonanalytic cutoff}
A compactly supported smooth cutoff cannot be treated as an analytic
function. We give the error explicitly. If $\beta$ is real with
$|\beta|\le h$, set $D_\beta=\beta\cdot\nabla_q$ and
\begin{equation}
 \chi_K(q+is\beta)=\sum_{j=0}^K\frac{(is)^j}{j!}D_\beta^j\chi(q),
 \qquad
 (\partial_s-iD_\beta)\chi_K
 =-\frac{i(is)^K}{K!}D_\beta^{K+1}\chi.               \label{c:eq:almostanalytic}
\end{equation}
Use $K=\lfloor[4(2^{16}/q_0)|\beta|]^{-1}\rfloor$; $\beta=0$ needs no
shift. Since $4\cdot2^{16}/q_0=2^{158}$ and $|\beta|\le h$, we have
$K\ge\lfloor L/2^{158}\rfloor\ge1$, because $\Ls\ge2^{158}$.
Lemma~\ref{c:lem:gevrey} and $K!\le K^K$ bound the last expression by
$\Cs e^{-\cs L}$. Its support is in the cutoff transition, where the
fiber inverse is uniformly bounded independently of $h$, including
through the limiting contour used below. The sum $\chi_K$ and each
of its first 32 derivatives are bounded by the fixed derivative
majorants in Section~\ref{c:app:ledger}.

Throughout, Bloch integrals carry the phase $e^{iq\cdot R}$, and the
contour is shifted by $+is\beta$ with $\beta\cdot R\ge0$.
For an analytic operator-valued integrand $Y$, differentiate
\[
 \mathcal J(s)=\int e^{i(q+is\beta)\cdot R}
       Y(q+is\beta)\chi_K(q+is\beta)\dd q.
\]
Holomorphy turns the derivative of the exponential and $Y$ into a
$q$-derivative. Integration by parts leaves exactly the last expression
in~\eqref{c:eq:almostanalytic}. There is no boundary term because all
cutoff derivatives vanish at the chart boundary. If $\beta\cdot R\ge0$,
the exponential has modulus at most one along the movement. Hence the
original integral differs from the shifted one by at most
$\Cs e^{-\cs L}$ times the explicitly bounded source or cell factors.
This is the only error introduced by moving the nonanalytic cutoff.

\subsection{The upper bound}
For real $q$ in the chart and
$|H^{1/2}\beta|\le .99h$, the even Taylor expansion gives
\begin{equation}
 \operatorname{Re}(\lambda(q+i\beta)-z_h-\sigma)
 \ge 2^{-8}(h^2+q^THq).                              \label{c:eq:tube}
\end{equation}
To verify it, the quadratic part is
$h^2+q^THq-\beta^TH\beta$ up to $2h^4$.
The quartic remainder and its mixed terms are absorbed using
$M_4(4q_0)^2<10^{-6}/\varrho$ and $h\le\Ls^{-1}$.
The other fiber directions remain separated by Lemma~\ref{c:lem:edge}.
The inverse on this tube is therefore bounded by a fixed constant times
$L^2$.

For a displacement $R$ choose
$\beta=.99h H^{-1}R/|H^{-1/2}R|$. The preceding contour calculation,
applied between two period cells, gives an upper bound
$\Cs L^2e^{-.99h|H^{-1/2}R|}+\Cs e^{-\cs L}$.
Each radius-$L$ support meets at most a fixed multiple of $L^3$ cells.
The row and column Schur bounds for the matrix of cell-block norms,
and the bounded change of phase across a support, give
\[
 \norm{Q_hR_0\mathcal T_RQ_h^*}
 \le\Cs L^5\big(e^{-.99h|H^{-1/2}R|}+e^{-\cs L}\big).
\]
For $R=2LdB\nu$,~\eqref{c:eq:Bmetric} and
$|\nu|_2\ge|\nu|_1/\sqrt3$ show that the exponent per $d|\nu|_1$ is at least
$2(.99)^2/\sqrt3>\log3$. This proves~\eqref{c:eq:blockdecay}.
The elementary comparison follows, for example, from
$\sqrt3<7/4$ and $\log3<11/10$, both certified by rational series.

\subsection{The lower bound: pole and transverse Gaussian}
For direction $e_j$, let $\beta_{j,h}>0$ be the unique small solution of
\begin{equation}
 E_c-\lambda(i\beta_{j,h}e_j)=E_c-z_h.
                                                               \label{c:eq:betaroot}
\end{equation}
The even expansion, positivity of $H$, and the derivative of its
quadratic term prove existence and uniqueness by monotonicity. They also
give
\begin{equation}
 .99/\sqrt{H_{jj}}\le\beta_{j,h}/h\le1.01/\sqrt{H_{jj}},
 \quad I_{j,h}:=2LB_{jj}\beta_{j,h}
 \le2(1.01)^2<\log8-1/50.                            \label{c:eq:action}
\end{equation}
Coordinate reflection symmetry makes the transverse gradient zero at
this point. For $0\le s<1$ the contour shifted by
$is\beta_{j,h}e_j$ is nonsingular. At $s=1$, its only selected-band
singularity is $q=0$. The real part of the denominator is at least a
fixed positive multiple of $|q|^2$ there, so the inverse is bounded by
$C(1+|q|^{-2})$, integrable in three dimensions. Dominated convergence
justifies the limiting contour. All errors from
\eqref{c:eq:almostanalytic} remain supported away from this singularity.

Put $q=hk$, $b_{j,h}=\beta_{j,h}/h$, and
\[
 D_h(k)=h^{-2}(\lambda(hk)-z_h).
\]
After the shift the selected contribution has the form
\begin{equation}
 L^2e^{-I_{j,h}d}
 \int e^{2iB_{jj}dk_j}
 \frac{\mathcal N_h(k+ib_{j,h}e_j)}
 {D_h(k+ib_{j,h}e_j)}\,
 \chi_K(hk+i\beta_{j,h}e_j)\dd k.                   \label{c:eq:scaledgreen}
\end{equation}
The normalization of $\mathcal N_h$ already includes $(2\pi)^{-3}$.
The other bands have the same exponential prefactor and a smooth
integrand. The cutoff movement adds at most $\Cs L^5e^{-\cs L}$.

Write $k=k_\parallel e_j+k_\perp$ with $k_\perp\perp e_j$. Near zero the denominator in
\eqref{c:eq:scaledgreen} has
\[
 D(0,0)=0,\quad D_\parallel(0,0)=ic_{j,h},\quad D_\perp(0,0)=0,
 \qquad c_{j,h}>0,
\]
with real positive transverse Hessian $Q_{j,h}$. The explicit bounds
above give
\begin{equation}
 1<c_{j,h}<2^{112},\qquad
 2^{98}I<Q_{j,h}<2^{110}I.                            \label{c:eq:polebounds}
\end{equation}
On a complex radius-one neighborhood after scaling, the analytic
functions and their derivatives have bounds independent of $h$;
$h\le\Ls^{-1}$ keeps this neighborhood far inside the Bloch chart.
The elementary analytic implicit-function construction, by contraction
in the $k_\parallel$ variable, gives the unique pole
\begin{equation}
 k_\parallel=p_h(k_\perp)=\frac{i}{2c_{j,h}}k_\perp^TQ_{j,h}k_\perp+O(|k_\perp|^3),
 \qquad \operatorname{Im}p_h(k_\perp)\ge2^{-256}|k_\perp|^2
                                                               \label{c:eq:pole}
\end{equation}
on $|k_\perp|\le2^{-4096}$. The explicit contraction and derivative bounds
are recorded in Section~\ref{c:app:ledger}. In particular this is a
uniform neighborhood, not one chosen separately for each input.

Factor $D(k_\parallel,k_\perp)=(k_\parallel-p_h(k_\perp))d_1(k_\parallel,k_\perp)$ and divide the analytic numerator:
\[
 \frac{\mathcal N_h(k_\parallel e_j+k_\perp+ib_{j,h}e_j)}{D(k_\parallel,k_\perp)}
   =\frac{\mathcal A_h(k_\perp)}{k_\parallel-p_h(k_\perp)}+\mathcal W_h(k_\parallel,k_\perp),
 \quad \mathcal A_h(0)=\frac{\mathcal N_h(ib_{j,h}e_j)}{ic_{j,h}}.
\]
Choose smooth cutoffs equal to one on a smaller fixed neighborhood,
with all poles over the transverse support inside the core of the
longitudinal cutoff. For $T=2B_{jj}d>0$ and $k_\perp\ne0$,
\[
 \int_{\R}\frac{e^{iTu}}{k_\parallel-p_h(k_\perp)}\dd k_\parallel=2\pi i e^{iTp_h(k_\perp)}.
\]
This identity follows by closing the contour in the upper half-plane;
alternatively integrate on a finite interval first. The oscillatory
tails are bounded by $C/(TR)$ by integration by parts and then tend
to zero. Replacing the full line by the longitudinal cutoff has error
$O(T^{-16})$, since the removed region stays a fixed distance from the
poles and its positive-order derivatives are integrable.

Now put $k_\perp=T^{-1/2}y$. Taylor's formula and~\eqref{c:eq:pole} give
\[
 \int \mathcal A_h(k_\perp)e^{iTp_h(k_\perp)}\dd k_\perp
 =\mathcal A_h(0)\frac{2\pi c_{j,h}}{T\sqrt{\det Q_{j,h}}}
    +O(T^{-3/2}).
\]
Indeed the amplitude error is bounded by $C|k_\perp|$ and the exponent error
by $CT|k_\perp|^3$, both against $e^{-2^{-257}T|k_\perp|^2}$. Their radial
integrals are bounded by fixed constants times $T^{-3/2}$.
The omitted Gaussian tails have the same bound after increasing the
fixed constant. Thus the leading coefficient in~\eqref{c:eq:scaledgreen}
is
\begin{equation}
 b^0_{j,h}=\frac{4\pi^2\mathcal N_h(ib_{j,h}e_j)}
 {2B_{jj}\sqrt{\det Q_{j,h}}},\qquad
 2^{-512}\le b^0_{j,h}\le\Cs.                         \label{c:eq:leadingcoefficient}
\end{equation}
The lower bound follows from~\eqref{c:eq:numeratorlower},
$B_{jj}<2^{53}$, and~\eqref{c:eq:polebounds}. In particular the leading
term cannot cancel between polarizations.

We also need a uniform bound on the remainder outside the pole chart,
whose scaled domain grows like $1/h$. For $|k|$ outside a fixed
neighborhood, the selected denominator and its derivatives satisfy
\[
 |D_h(k+ib_{j,h}e_j)|\ge c(1+|k|^2),\qquad
 |\partial^K D_h(k+ib_{j,h}e_j)|
 \le C_K(1+|k|)^{2-K}.
\]
For orders $K\ge3$ use $h^{K-2}$ from differentiation and
$h\le4q_0/|k|$ on the scaled chart. The inverse-derivative recurrence
then gives the corresponding order $-2$ symbol bounds. Each fixed
numerator derivative is bounded by Lemma~\ref{c:lem:numerator}. The
localized regular integrand consists of the regular quotient $\mathcal W_h$ times
the pole-chart cutoffs, the selected quotient outside the pole chart,
including the transition shell of $\chi_K$, and the other bands. Hence the
$L^1$ norm of the sixteenth longitudinal derivative of the localized
regular integrand is at most $\Cs/h$. The integral of
$(1+|k|)^{-2}$ over a radius-$4q_0/h$ ball is $O(1/h)$.
For the other bands the scaled inverse includes a factor $h^2$ and
its integral over that ball is also $O(1/h)$. Derivatives of the
shifted cutoff are bounded uniformly, as shown after
\eqref{c:eq:almostanalytic}. Sixteen integrations by parts therefore
bound this entire regular contribution by $\Cs Ld^{-16}$ inside the
brackets of the following estimate:
\begin{equation}
 \left|h_j(d)-L^2e^{-I_{j,h}d}\frac{b^0_{j,h}}d\right|
 \le\Cs L^2e^{-I_{j,h}d}
       (d^{-3/2}+Ld^{-16})+\Cs L^5e^{-\cs L}.          \label{c:eq:greenremainder}
\end{equation}
This calculation includes the full regular spectral complement, not
only the selected-band numerator. Here, as throughout, $\Cs$ is the
fixed number in~\eqref{c:eq:masterconstants}.

For $d\ge\Cs^8$ and $d^{15}\ge\Cs^4L$, the first two errors in
\eqref{c:eq:greenremainder} are less than one quarter of the main term.
Equation~\eqref{c:eq:action} gives a strict exponential margin; for these
$d$, $e^{d/50}/d$ is larger than $2^{1024}$. The last error is absorbed
by~\eqref{c:eq:couplingtail}. This proves~\eqref{c:eq:couplinglower} and
completes Lemma~\ref{c:lem:interactions}.

\begin{lemma}[Off-support and periodization estimates]
\label{c:lem:cellbounds}
Let $\Box_a=2a+[0,2)^3$ be period cells. For $0\le\sigma\le h^4/64$,
let $\mathscr R_\sigma$ denote the local fiber inverse on the support
of $\chi$, before multiplication by $\chi$, and define
\begin{equation}
 R_\sigma^{(j)}=
 \chi\mathscr R_\sigma(\partial_j\varepsilon_{\rm per})
       \mathscr R_\sigma.
                                                               \label{c:eq:Rderivative}
\end{equation}
For $\mathcal R=R_\sigma$ or $\mathcal R=R_\sigma^{(j)}$,
\begin{align}
 \norm{\one_{\Box_a}\mathcal R\one_{\Box_b}}
 &\le\Cs L^4\big(e^{-\cs\dist(\Box_a,\Box_b)/L}+e^{-\cs L}\big),
                                                               \label{c:eq:cellspatial}\\
 \norm{\one_{\Box_a}\mathcal R\one_{\Box_b}}
 &\le\Cs L^4 e^{-\cs\sqrt{|a-b|/L}}.                 \label{c:eq:cellgevrey}
\end{align}
For $f_{\rm a}\in H^1$,
\begin{equation}
 \partial_j(R_\sigma f_{\rm a})=R_\sigma\partial_jf_{\rm a}+
                         R_\sigma^{(j)}f_{\rm a}.            \label{c:eq:derivativeidentity}
\end{equation}
The auxiliary source $f_{\rm a}$ of Section~\ref{c:sec:assembly} belongs to $H^1$
by Lemma~\ref{c:lem:auxrecovery}.
\end{lemma}
\begin{proof}
The first estimate repeats the nonsingular tube argument. A product of
two local inverses has norm at most a fixed constant times $L^4$ on
a slightly smaller tube. Equation~\eqref{c:eq:Bmetric} is not needed:
the smallest Euclidean decay rate is at least a fixed multiple of
$h/\sqrt{\norm A}$, larger than $\cs h$.

For the second estimate, the cell-fiber integrand is analytic in a
complex strip of width at least $2^{-256}h$ over the real support of
the cutoff. Cauchy's estimate, the Gevrey bounds on $\chi$, and the
bounded cell gauge factors give derivatives of order $K$ at most
\[
 \Cs L^4(2^{512}L)^K(K!)^2.
\]
Integrating by parts in the cell Fourier coefficient and optimizing
$K$ as in Lemma~\ref{c:lem:gevrey} proves~\eqref{c:eq:cellgevrey}, with
the smaller constant $\cs$ in the exponent. This summable estimate,
rather than the nondecaying error term in~\eqref{c:eq:cellspatial}, is
used for periodization.

For~\eqref{c:eq:derivativeidentity}, translate the periodic material and
use the local resolvent identity, then differentiate at zero. The
cutoff commutes with translations. The formula contains
one factor $\chi$, not two, since $\chi$ is not a projection. The local
inverse without $\chi$ in~\eqref{c:eq:Rderivative} is essential for
this identity. The norm bounds justify the difference-quotient limit.
No $H^1$ estimate on individual cell restrictions or artificial cell
traces is used.
\end{proof}

\section{Leakage from the auxiliary channel}
\label{c:sec:leakage}

The estimates in this section are the replacement for a complete bulk
band gap. They control the actual error produced by the cutoff in
\eqref{c:eq:Ridentities}.

\begin{lemma}[A nested momentum-cutoff estimate]
\label{c:lem:leakage}
Let $r_i\in2\Z^3$, $1\le i\le J$, and suppose the balls
$\mathbb B_L+r_i$ are disjoint. Put
\[
 W(r)=\sum_i v_iw(h(r-r_i)),\qquad
 0\le v_i\le2^{130}h^2,
 \qquad a_i=\mathcal T_{r_i}a_h.
\]
If $f_* =\sum_i c_i a_i$ and
\begin{equation}
 f_{\rm a}=f_*+WR_\sigma f_{\rm a},                                  \label{c:eq:tequation}
\end{equation}
then
\begin{equation}
 \norm{(1-\chi)f_{\rm a}}\le
 \Cs(1+J)e^{-\cs L^{1/3}}
       \left(\norm{f_{\rm a}}+\sum_i|c_i|\right).             \label{c:eq:leakage}
\end{equation}
No support or polarization condition is imposed on $f_{\rm a}$ in this lemma.
\end{lemma}
\begin{proof}
Let $P_{\ge s}$ be the Bloch-momentum projection onto torus distance at least
$s$ from the corner. On the support of $\chi$ with $|q|\ge q_0/4$,
\eqref{c:eq:edgebounds} and~\eqref{c:eq:quartic} give
\begin{equation}
 \norm{P_{\ge s}R_\sigma}\le2^{192}\quad(s\ge q_0/4).
                                                               \label{c:eq:highqbound}
\end{equation}
For example the quadratic denominator at $q_0/4$ is at least
$2^{-187}$, and the other bands are separated. Disjointness gives
$\norm W\le2^{130}h^2$, so their product is less than $1/4$.

If the two Bloch-momentum sets are separated by $\Delta_q$, multiplication
by $W$ between them has norm at most
\begin{equation}
 2^{100}J\max_i v_i\,e^{-\sqrt{\Delta_q/h}/512}.         \label{c:eq:fouriercross}
\end{equation}
Indeed its physical Fourier convolution kernel is the sum of translated
copies of $v_ih^{-3}(\mathcal Fw)(p/h)$. Frequencies in the two sets
differ in physical momentum by at least their reciprocal-torus
distance $\Delta_q$. Young's inequality and~\eqref{c:eq:gevreytail}
prove~\eqref{c:eq:fouriercross}. Possible changes of reciprocal alias
do not invalidate this lower bound on physical momentum difference.

Apply $P_{\ge s}$ to~\eqref{c:eq:tequation}, split the input at $s-\Delta_q$,
and use that $R_\sigma$ commutes with these projections. The high-input
part is bounded by $\tfrac14\norm{P_{\ge s-\Delta_q}f_{\rm a}}$. The low-input
part is bounded by~\eqref{c:eq:fouriercross} times $4L^2\norm{f_{\rm a}}$.
The factor $h^2$ in $v_i$ cancels this $L^2$; the resulting constant
is at most $2^{232}J$.

First apply this argument to~\eqref{c:eq:source-mode-equation}. It holds
also for every translated $a_i$, because the centers are period
translations. Use $K=\lceil L^{1/3}\rceil$ nested cutoffs from
$q_0/4$ to $q_0/2$, with step $\Delta_q=q_0/(4K)$ and coefficient
$\mu_h^{-1}\le2^{128}h^2$. Iteration gives
\[
 \norm{P_{\ge q_0/2}a_i}\le
 4^{-K}\norm{a_i}+
 2^{233}e^{-\sqrt{q_0L/(4K)}/512}\norm{a_i}
 \le\Cs e^{-\cs L^{1/3}},
\]
since $\norm{a_i}\le1$. The last comparison uses
$K\le2L^{1/3}$ and $q_0=2^{-140}$; even the much larger exponent
constant $2^{-84}$ would suffice.

Now iterate the same inequality for $f_{\rm a}$ from $q_0/2$ to $q_0$.
Every source term is at most
$\Cs e^{-\cs L^{1/3}}\sum_i|c_i|$ by the preceding estimate.
The geometric sum of the contraction factors is less than two.
The initial remainder is $4^{-K}\norm{f_{\rm a}}$, and all low-input leakage
terms have the same exponential bound. Finally $1-\chi$ vanishes on
$|q|<q_0$, so its norm is bounded by that of $P_{\ge q_0}f_{\rm a}$.
This proves~\eqref{c:eq:leakage} for every component. It also covers
momentum outside the chart, where $R_\sigma$ is zero.
\end{proof}

For $E_{\rm a}=R_\sigma f_{\rm a}$ satisfying~\eqref{c:eq:tequation}, the exact
physical residual in the infinite medium is therefore
\begin{equation}
 (A_{\rm per}-z_h-\sigma-W)E_{\rm a}-f_*
       =-(1-\chi)f_{\rm a}.                                 \label{c:eq:physicalleakidentity}
\end{equation}
The construction below makes the right side smaller than the required
inverse-polynomial budget. The cutoff is not being inserted into the
physical equation.

\section{Finite algebra and assembly of the auxiliary response}
\label{c:sec:assembly}

\subsection{The counted logical network}\label{c:sec:network}
We use the logical interface of Lemma~\ref{c:lem:logical}.


Thus $1\le\norm g\le\sqrt m$.
Double the logical lattice. Put a terminal at each doubled vertex and
one mediator at the midpoint of each actual edge. If $J$ is the total
number of these sites, then $J\le2^{19}\ell_{\rm p}^3$. Multiply their coordinates
by $2LdB$. All centers are in $2\Z^3$. At placement-lattice distance
one the only occupied pairs are the intended terminal--mediator
incidences: terminals have all coordinates even, and a midpoint has
exactly one odd coordinate. Every other occupied pair is at distance
at least two. This includes the two vertices of a parity detour in
Section~\ref{d:app:routing}, which are lattice-adjacent but not joined
by a path edge: after doubling they lie at distance two with no mediator
between them. The physical balls are disjoint.

\subsection{Explicit separation and diagonal tuning}
For the moment suppose $L,d$ satisfy the tests below; their terminating
selection is proved in Section~\ref{c:sec:parameters}. Put
\begin{equation}
 \zeta=3^{-d},\qquad
 \alpha=(c_u/4)L^2 8^{-d}=2^{-1026}L^2 8^{-d}.         \label{c:eq:alpha}
\end{equation}
Require
\begin{align}
 d&\ge\ds, &d^{15}&\ge\Cs^4L, &\alpha&\le1,\\
 \Cs^4L^{12}\zeta&\le\ee/2^{24},
 &\Cs^4L^{12}\zeta^2&\le\alpha\ee/(2^{24}\overline K),          \label{c:eq:dtests}\\
 J\Cs L^5e^{-\cs L}&\le\zeta^2/4.                   \label{c:eq:Ltailtest}
\end{align}
The last two inequalities imply~\eqref{c:eq:couplingtail}.
Calibrate the three real numbers $-h_j(d)$ by dyadic $\beta_j$ with
absolute error at most $\zeta^2/16$. Lemma~\ref{c:lem:interactions} gives
$|\beta_j|\ge3\alpha$.

Order terminals first and mediators second. Define the target matrix
$\Theta$ by incidences $\beta_j$ for a mediator in direction $j$ and
\begin{equation}
 \Theta_{ii}=\alpha(\Gamma_{ii}-\deg i),\qquad
 \Theta_{ee}=-\beta_j^2/\alpha.                            \label{c:eq:T}
\end{equation}
Its mediator Schur complement is exactly $\alpha\Gamma$. If $\Theta_{\rm tm},\Theta_{\rm mm}$
are its incidence and mediator blocks, then
\begin{equation}
 \norm{\Theta_{\rm mm}^{-1}}\le\alpha^{-1},\quad
 \norm{\Theta_{\rm tm}\Theta_{\rm mm}^{-1}}\le\sqrt6,\quad
 \norm{\Theta^{-1}}\le16\overline K/\alpha.                         \label{c:eq:Tinverse}
\end{equation}
Indeed each mediator contributes
$-\beta_j(-\beta_j^2/\alpha)^{-1}\beta_j=\alpha$ to both endpoint
diagonals and to their mutual edge. The endpoint diagonal correction
cancels the degree term in~\eqref{c:eq:T}. Each entry of
$\Theta_{\rm tm}\Theta_{\rm mm}^{-1}$ has magnitude at most $1/3$; its row and column degrees
are at most three and two. Triangular block factorization proves
\eqref{c:eq:Tinverse}. For any residual vector,
\begin{equation}
 \norm{\Gamma z_{\rm t}-g}\le(4/\alpha)
             \norm{\Theta z-\alpha g_\Theta}.                 \label{c:eq:mediatorres}
\end{equation}

Choose positive dyadic $v_i$ so that
\begin{equation}
 |v_i^{-1}-(\mu_h+\Theta_{ii})|\le\zeta^2.                \label{c:eq:vtuning}
\end{equation}
For terminals, $|\Theta_{ii}|\le6\alpha\le6$ and $L\ge\Ls$ give
$|\Theta_{ii}|\le2^{-196}L^2$. For mediators, $|\beta_j|\le2\Cs L^5\zeta$
and the second smallness test give the same bound.
Thus the tuning interval following Lemma~\ref{c:lem:numerator} applies,
and $0<v_i\le2^{130}L^{-2}<1$.
To implement~\eqref{c:eq:vtuning}, enclose $\mu_h$ to $\zeta^2/16$,
form the rational target, and round its reciprocal on a dyadic grid
of spacing at most $\zeta^2/(16L^4)$. The reciprocal derivative is
at most $4L^4$ on this interval. This is a finite precision rule,
not an exact-real tuning instruction.

Let $\tau$ be the largest power of two not exceeding
\begin{equation}
 \frac{\zeta^2}{2^{24}L^4}.                           \label{c:eq:tau}
\end{equation}
All $0\le\sigma\le\tau$ satisfy~\eqref{c:eq:Ridentities}; in particular
$\norm{R_\sigma-R_0}\le\zeta^2/2^{20}$.
Let $Q$ collect the translated $Q_h$ blocks, let $V_{\rm d}=\diag(v_iI)$,
and set
\begin{equation}
 W=Q^*V_{\rm d}Q,\qquad
 \MM=V_{\rm d}^{-1}-QR_\sigma Q^*.                          \label{c:eq:Mdef}
\end{equation}
Here the defect direct sum is infinite dimensional, and $\norm Q\le1$.

\begin{lemma}[Stable auxiliary assembly]
\label{c:lem:assembly}
For every $0\le\sigma\le\tau$, $\MM$ is invertible and
$\norm{\MM^{-1}}\le256\overline K/\alpha$. If $U$ inserts $u_h$ at every site,
then its bright Schur block has the form
\begin{equation}
 S_b=\Theta+\Delta,\qquad
 \norm\Delta\le\Cs^4L^{12}\zeta^2,
 \qquad \norm{S_b^{-1}}\le32\overline K/\alpha.                \label{c:eq:bright}
\end{equation}
All complementary modes and all unwanted interactions are included.
\end{lemma}
The parity layout of Section~\ref{c:sec:network} is essential here:
it places every unwanted occupied pair at $|\nu|_1\ge2$, so that all
unwanted couplings enter the non-neighbor sum of order $\zeta^2$ below.
Lemma~\ref{c:lem:interactions} alone does not exclude unwanted pairs at
$|\nu|_1=1$, whose coupling would be of the useful size.
\begin{proof}
For $0<\zeta\le1/2$, the elementary identity
$\sum_{\nu\in\Z^3}\zeta^{|\nu|_1}=((1+\zeta)/(1-\zeta))^3$ gives
\[
 \sum_{\nu\ne0}\zeta^{|\nu|_1}\le64\zeta,
 \qquad\sum_{|\nu|_1\ge2}\zeta^{|\nu|_1}\le256\zeta^2.
\]
The matrix of block norms bounds the complete operator by its row and
column sums. The geometric part of~\eqref{c:eq:blockdecay} therefore
has norm at most $64\Cs L^5\zeta$, and its non-neighbor part at most
$256\Cs L^5\zeta^2$. The nondecaying contour error is summed only
over the finite collection:~\eqref{c:eq:Ltailtest} bounds it by
$\zeta^2/4$. It has not been summed over an infinite lattice.
The uniform detuning changes the whole compressed union by at most
$\zeta^2/2^{20}$, independently of $J$.

Write the bright/complement decomposition of $\MM$ as
$\left(\begin{smallmatrix}\MM_{bb}&\MM_{bc}\\\MM_{bc}^*&\MM_{cc}\end{smallmatrix}\right)$.
At infinite separation its complementary block is the direct sum of
$\mu_hI-S_h$ on $u_h^\perp$, with gap $2^{-192}L^2$.
On-site scalar shifts have no mixed block. The displayed row sums and
\eqref{c:eq:dtests} give
\[
 \norm{\MM_{cc}^{-1}}\le2^{194}L^{-2},\qquad
 \norm{\MM_{bc}}\le128\Cs L^5\zeta,\qquad
 1+\norm{\MM_{bc}}\norm{\MM_{cc}^{-1}}\le2.
\]
This is an inverse-perturbation argument, valid even though the
complement may be indefinite. Direct bright errors consist of diagonal
rounding, incidence calibration, detuning, and the non-neighbor sum.
Adding $\MM_{bc}\MM_{cc}^{-1}\MM_{bc}^*$ bounds their sum by
$\Cs^4L^{12}\zeta^2$. Equations~\eqref{c:eq:Tinverse} and
\eqref{c:eq:dtests} make $\norm{\Theta^{-1}\Delta}<1/2$, proving the inverse
bound for $S_b$. A second triangular factorization, with factors of
norm at most two, gives $\norm{\MM^{-1}}\le256\overline K/\alpha$.
\end{proof}

\subsection{\texorpdfstring{A local $L^2$ source and its auxiliary logical response}{A local L2 source and its auxiliary logical response}}
Put $g_\Theta=(g,0)$ and define the ideal source
\begin{equation}
 f_*=(\alpha/\mu_h)Q^*Ug_\Theta.                          \label{c:eq:idealsource}
\end{equation}
It is a compactly supported $H^1$ function, hence a genuine $L^2$ source.
No unbounded operator is applied in its definition. Define, only for
analysis,
\begin{align}
 w_d&=\MM^{-1}QR_\sigma f_*,\\
 E_{\rm a}&=R_\sigma(f_*+Q^*w_d),\qquad
 f_{\rm a}=f_*+Q^*w_d.                                      \label{c:eq:auxsolution}
\end{align}
Then $w_d=V_{\rm d}QE_{\rm a}$ and $f_{\rm a}=f_*+WE_{\rm a}$, so
\eqref{c:eq:tequation} holds. The encoder never computes this inverse or
this field.

\begin{lemma}[Auxiliary recovery and magnitude bounds]
\label{c:lem:auxrecovery}
Let $z=U^*V_{\rm d}QE_{\rm a}$, i.e., $z_i=v_i\ip{E_{\rm a}}{a_i}$ are the
volume moments~\eqref{c:eq:idealreadout} of $E_{\rm a}$. Then
\begin{equation}
 \norm{\Gamma z_{\rm t}-g}\le\ee\norm g/1024.                  \label{c:eq:auxlogical}
\end{equation}
Also
\begin{align}
 2^{-256}\alpha L^{-2}\norm g&\le\norm{f_*}\le\norm g,\\
 \norm{f_{\rm a}}&\le\Cs \overline K\norm g,
 &\norm{E_{\rm a}}&\le\Cs \overline KL^2\norm g,\\
 \norm{f_{\rm a}}_{H^1}&\le\Cs^4\overline K^2\alpha^{-1}L^4\norm g,
 &\norm{E_{\rm a}}_{H^1}&\le\Cs^5\overline K^2\alpha^{-1}L^6\norm g.
                                                               \label{c:eq:auxnorms}
\end{align}
\end{lemma}
\begin{proof}
On each diagonal block $Q_hR_0Q_h^*u_h=\mu_hu_h$.
The off-site row sums and the detuning bound therefore give
\[
 QR_\sigma f_*=\alpha Ug_\Theta+\rho_s,
 \qquad\norm{\rho_s}\le\alpha\ee\norm g/2^{24}.
\]
Here the factor $\mu_h^{-1}\le2^{128}L^{-2}$ is absorbed by the first
smallness test in~\eqref{c:eq:dtests}. Eliminating the complement
multiplies this error by at most two. Thus
$\norm{S_bz-\alpha g_\Theta}\le\alpha\ee\norm g/2^{23}$ and
$\norm z\le33\overline K\norm g$. Adding $\Delta z$, using the second test, and
applying~\eqref{c:eq:mediatorres} proves~\eqref{c:eq:auxlogical} with
room to spare.

Since $\mu_h=\ip{a_h}{R_0a_h}\le4L^2\norm{a_h}^2$,
$\norm{a_h}\ge2^{-65}$. Disjoint supports,~\eqref{c:eq:mugap}, and
\eqref{c:eq:alpha} give the source bounds. From
\eqref{c:eq:Ridentities} and Lemma~\ref{c:lem:assembly},
\[
 \norm{QR_\sigma f_*}\le2^{130}\alpha\norm g,
 \quad\norm{w_d}\le2^{138}\overline K\norm g,
 \quad\norm{E_{\rm a}}\le2^{141}\overline KL^2\norm g,
\]
which are smaller than the stated master bounds.

For regularity, the exact identity
\begin{equation}
 (I-WR_\sigma)^{-1}=I+Q^*\MM^{-1}QR_\sigma            \label{c:eq:t-inverse}
\end{equation}
has norm at most $1+1024\overline KL^2/\alpha$. It is checked by direct
multiplication. Translate~\eqref{c:eq:tequation}. The source derivative
is bounded using~\eqref{c:eq:modeH1}. Disjointness gives
$\norm{\nabla W}_\infty\le\Cs h^3$. The resolvent $R_\sigma$ commutes
only with period translations; conjugation by a translation $\mathcal T_a$ moves
$\varepsilon_{\rm per}$, and the resolvent identity used
for~\eqref{c:eq:modeH1} gives
$\norm{\mathcal T_aR_\sigma\mathcal T_a^{-1}-R_\sigma}\le\Cs L^4|a|$.
Subtracting the translated equation and applying~\eqref{c:eq:t-inverse}
therefore bounds its difference quotient by
$\Cs^4\overline K^2\alpha^{-1}L^4G$. The Fourier difference-quotient argument
proves $f_{\rm a}\in H^1$. Equation~\eqref{c:eq:derivativeidentity} then gives
the last bound in~\eqref{c:eq:auxnorms}. Both $f_{\rm a}$ and its weak first
derivatives are supported in the union of the well balls.
\end{proof}

\section{Finite PEC boundaries and the actual material}
\label{c:sec:boundary}

Choose an even integer-sided box containing every well ball at distance
at least $\ell_{\rm buf}+4$ from its boundary. Translate the whole layout by an even
integer vector if needed. Let $\vartheta_\Omega$ be a real $C^2$ cutoff,
equal to one at distance at least two from the boundary and zero at
distance at most one. It can be made as a product of the quintic ramp
$10s^3-15s^4+6s^5$ on $[0,1]$ and its reflected copies; its first two
derivatives are bounded by $2^{20}$. Put
\begin{equation}
 E_{\rm ref}=\vartheta_\Omega E_{\rm a}.                \label{c:eq:cutfield}
\end{equation}
It lies in $D(\Cc)$ for the finite PEC box. Indeed $E_{\rm a}\in H^1$
by~\eqref{c:eq:auxnorms},
its curl--curl belongs to $L^2$ by~\eqref{c:eq:Ridentities}, and the
commutator with a $C^2$ cutoff is first order in $E_{\rm a}$.
The resulting field has compact support inside the box.

Define the explicit majorant
\begin{equation}
 \Xi=\Cs^{32}(1+N_\Omega)^2(1+J+m)^2
                  (1+\overline K)^4(1+\alpha^{-1})^2L^{100}.   \label{c:eq:Pglobal}
\end{equation}
\begin{lemma}[A physical finite-box approximate solution]
\label{c:lem:cutresidual}
For the ideal scalar material
$\varepsilon_{\rm id}=\varepsilon_{\rm per}+z_h+\sigma+W$,
\begin{equation}
 \norm{(\Cc-\varepsilon_{\rm id})E_{\rm ref}-f_*}
 \le\Xi\big(e^{-\cs L^{1/3}}+e^{-\cs \ell_{\rm buf}/L}\big)\norm g.
                                                               \label{c:eq:cutresidual}
\end{equation}
Furthermore $\norm{E_{\rm ref}}\le\Cs \overline KL^2\norm g$, and all ideal readout
moments on the wells are unchanged by the cutoff.
\end{lemma}
\begin{proof}
Equation~\eqref{c:eq:physicalleakidentity} gives the bulk error. In
Lemma~\ref{c:lem:leakage}, $\sum_i|c_i|\le
(\alpha/\mu_h)\sqrt m\,\norm g$, so~\eqref{c:eq:auxnorms} bounds it by the
first term of~\eqref{c:eq:cutresidual}.

The cutoff commutator is
\[
 [\curl\curl,\vartheta_\Omega]E
   =\nabla\vartheta_\Omega\times\curl E
       +\curl(\nabla\vartheta_\Omega\times E).
\]
Its norm is bounded by $2^{24}$ times the $H^1$ norm of $E$ on the
boundary strip. Apply~\eqref{c:eq:cellspatial} to the compactly supported
$f_{\rm a}$ and $\partial_jf_{\rm a}$, and use~\eqref{c:eq:derivativeidentity}.
The matrix of cell-block norms between that support and the boundary
strip has norm at most its entry bound times the geometric mean of
the two cell counts, at most a fixed multiple of $N_\Omega$.
Their cell separation is at least $\ell_{\rm buf}$ up to the already reserved fixed
margin. Thus the boundary $H^1$ norm is bounded by
\[
 \Cs^8 N_\Omega \overline K^2\alpha^{-1}L^{10}
        (e^{-\cs \ell_{\rm buf}/L}+e^{-\cs L})\norm g.
\]
The constant and powers are smaller than~\eqref{c:eq:Pglobal}.
Since $L\ge1$, $e^{-\cs L}\le e^{-\cs L^{1/3}}$.
This proves the residual bound. The cutoff equals one on the wells,
so the other assertions follow immediately.
\end{proof}

The actual output material uses a finite radial polynomial $\widetilde p_{\rm r}(t)$
with dyadic coefficients, approximating $e^{-1/[2(1-t)]}$ on
$0\le t\le1$, with $\widetilde p_{\rm r}(1)=\widetilde p_{\rm r}'(1)=0$. Outside each ball it is zero.
Set
\begin{equation}
 \overline\varepsilon_0(r)=7+z_h+\varrho V_0(r)
       +\sum_i v_i \widetilde p_{\rm r}(|r-r_i|^2/L^2)^2\one_{|r-r_i|<L},
 \qquad \overline\varepsilon=\overline\varepsilon_0+\sigma.
                                                               \label{c:eq:actualmaterial}
\end{equation}
Section~\ref{c:sec:representation} constructs this polynomial with a
prescribed uniform error and proves its cost. The polynomial square is
nonnegative. The background is between six and eight; the well
increments are less than one, and all approximation budgets are less
than $1/16$. Thus $1\le\overline\varepsilon\le25$ on every draw
(in fact it stays between six and ten).

Let
\begin{equation}
 e_0=\frac{\ee}{2^{24}J_\Gamma L_b},\qquad J_\Gamma=8.
                                                               \label{c:eq:e0}
\end{equation}
Here $J_\Gamma=8$ is an upper bound for $\norm\Gamma\le7$, which
multiplies the moment errors in~\eqref{c:eq:final-logical}.
Choose the material approximation so that
\begin{equation}
 \norm{\overline\varepsilon-
          (\varepsilon_{\rm per}+z_h+\sigma+W)}_\infty
       \le\frac{e_0}{4\Cs \overline KL^2}.                    \label{c:eq:materialprecision}
\end{equation}
Its additional residual on $E_{\rm ref}$ is at most $e_0\norm g/4$.
This bounded-multiplication comparison is why the output polynomial
need not retain the ideal well's Gevrey smoothness or exact Bloch
localization. It is not justified by qualitative convergence.
The random detuning of Lemma~\ref{c:lem:detuning} is applied to the
\emph{actual} material~\eqref{c:eq:actualmaterial}, after its finite
approximation has been fixed.

\section{Terminating parameter selection}
\label{c:sec:parameters}

All choices in this section use integers and rational comparisons. In
particular, the encoder does not need a real logarithm or a decision
about an exact transcendental equality. The logical parameters
$\overline K,\ee,J$ and the occupied sites have already been computed.
Write $\ell=1+\lceil\log_2(2+m)\rceil$.

For a candidate power of eight $L\ge\max(\Ls,\ell^6)$, starting with the
least such power and multiplying by eight after each failure, choose $d$ by
starting at one and doubling until the rational tests
\eqref{c:eq:dtests} hold. The tests include $d\ge\ds$,
$d^{15}\ge\Cs^4L$, and $\alpha\le1$; the definitions are
$\zeta=3^{-d}$ and $\alpha=2^{-1026}L^2 8^{-d}$.
Let $\tau$ be the largest integral power of two no greater than
$\zeta^2/(2^{24}L^4)$. This implies
\begin{equation}
 \frac{\zeta^2}{2^{25}L^4}<\tau\le
 \frac{\zeta^2}{2^{24}L^4}.                         \label{c:eq:tau-bracket}
\end{equation}
Every center is $2LdB$ times its doubled logical coordinate, followed
by one even integer translation. The extrema of these centers are
computed directly.

For $t_{\rm b}$ a power of two, starting at one, put $\ell_{\rm buf}=Lt_{\rm b}$. Choose the smallest
even integer side lengths containing the centers with margin at least
$L+\ell_{\rm buf}+4$ on all sides. Compute $N_\Omega,L_b$ by
\eqref{c:eq:count} and~\eqref{c:eq:detuning}, and compute $e_0$ and
$\Xi$ by~\eqref{c:eq:e0} and~\eqref{c:eq:Pglobal}. Double $t_{\rm b}$ until
\begin{equation}
 \Xi\,2^{-\lfloor\cs \ell_{\rm buf}/L\rfloor}\le e_0/4.
                                                               \label{c:eq:buffer-test}
\end{equation}
The dependence of $L_b$ and $\Xi$ on the candidate box is
included at every test. Only dimensions and counts, not the ambient
array, are written during this search.

After the buffer test succeeds, require also
\begin{align}
 \Xi\,2^{-\lfloor\cs L^{1/3}\rfloor}&\le e_0/4,
                                                               \label{c:eq:leak-test}\\
 J\Cs L^5\,2^{-\lfloor\cs L\rfloor}&\le\zeta^2/4.
                                                               \label{c:eq:tail-test}
\end{align}
If either test fails, replace $L$ by $8L$ and repeat. Because $L$ is a
power of eight, its cube root is an integer. The floors here are floors
of explicit rationals; they are ordinary integer divisions. The
inequality $e^{-s}\le2^{-\lfloor s\rfloor}$ for $s\ge0$ proves that
these rational tests imply all the corresponding exponential error
budgets. Test~\eqref{c:eq:tail-test} implies
\eqref{c:eq:couplingtail}: the second spacing test gives
$\zeta^2\le c_uL^2 8^{-d}$ with ample margin.

\begin{lemma}[Termination and size]
\label{c:lem:parameters}
The preceding procedure terminates. Its output satisfies
\begin{equation}
 L=O(\ell^6),\qquad d=O(\ell),\qquad \ell_{\rm buf}=O(\ell^7).
                                                               \label{c:eq:polylogscales}
\end{equation}
Its arithmetic cost before the final geometry is written is polynomial
in $\ell$, with an effectively specified fixed multiplicative constant.
All inverse scales in the construction are polynomial in $m$.
\end{lemma}
\begin{proof}
Here are explicit bounds, including the fixed constants. Put
\begin{equation}
 Z=\Cs^{128}(1+\lceil c_\kappa\rceil+\lceil c_\eta\rceil+
                         \lceil B_{\rm in}\rceil)^8,
 \qquad Z_L=8^{Z^8}.                                \label{c:eq:ZAL}
\end{equation}
The formulas in~\eqref{c:eq:logicalparameters}, using $n\le m$ and
$N_{\rm r}\le 8+(\lceil c_\kappa\rceil+\lceil c_\eta\rceil+2)\log_2m$, give
\[
 \log_2(2+\overline K+\ee^{-1}+J)\le Z^{1/4}\ell.
\]
Here the residual factor is
$T_{\rm res}=O(m^{7/2}\ell^{5/2})$ and
$\ee^{-1}=O(m^{\lceil C\rceil+11/2}\ell^{5/2})$.
The logarithms of the fixed coefficients and the displayed exponents are
bounded by $Z^{1/4}$ with the stated value of $Z$.
The power $Z^{1/4}$ only appears in this upper bound; the encoder never
computes it.

The first spacing inequality follows once a constant multiple of
$\log_2(\ee^{-1})+\log_2L+\log_2\Cs+1$ is at most $d$.
For the second, division by $\alpha$ leaves
\[
 2^{1050}\Cs^4 L^{10}\overline K\ee^{-1}(8/9)^d\le1.
\]
The elementary inequality $(9/8)^8>2$ supplies an entirely rational
sufficient test. The constraints $d\ge\ds$, $d^{15}\ge\Cs^4L$,
and $\alpha\le1$ have the same form or the explicitly displayed
fifteenth root. Doubling therefore stops with
\begin{equation}
 d\le Z(\ell+\log_2L+L^{1/15}).                    \label{c:eq:dcoarse}
\end{equation}
For example, to avoid the root in a sufficient candidate, the least
integer whose fifteenth power exceeds $\Cs^4L$ can be found by
integer doubling and bisection. Its value is at most
$2\Cs^{4/15}L^{1/15}$, which is covered by~\eqref{c:eq:dcoarse}.

The balanced logical box and $B_{jj}<2^{53}$ bound each physical side
by
\[
 2^{70}(1+\lceil c_\kappa\rceil+\lceil c_\eta\rceil)
 Ld\sqrt{m\ell}+2\ell_{\rm buf}+16.
\]
Together with~\eqref{c:eq:tau-bracket}, this bounds the logarithms of
$N_\Omega,L_b,e_0^{-1}$ and $\Xi$ by a fixed multiple of
\[
 \log_2\Cs+\ell+d+\log_2L+\log_2(1+t_{\rm b}).
\]
Substitution of the displayed formulas, allowing coefficient $Z^{1/4}$,
is a conservative bound for that fixed multiple and all additive
constants. Since $\cs=2^{-1024}$ and $Z\ge2^{2048}$, the first power
of two $t_{\rm b}$ above $Z^3(\ell+d+\log_2L)$ satisfies
\eqref{c:eq:buffer-test}. Indeed its linear left exponent
$\lfloor\cs t_{\rm b}\rfloor$ exceeds the stated logarithmic bound, using
$\log_2t_{\rm b}\le3\log_2Z+\log_2(\ell+d+\log_2L)+1$ and
$\log_2s\le s$ for $s\ge1$. Thus
\begin{equation}
 t_{\rm b}\le Z^4(\ell+d+\log_2L).                         \label{c:eq:tcoarse}
\end{equation}
After this buffer has been chosen, all logarithmic right-hand budgets
in~\eqref{c:eq:leak-test}--\eqref{c:eq:tail-test} are at most
\begin{equation}
 Z^6(\ell+\log_2L+L^{1/15}).                        \label{c:eq:budgetcoarse}
\end{equation}
This follows by inserting~\eqref{c:eq:dcoarse}--\eqref{c:eq:tcoarse};
$\log_2t_{\rm b}$ adds only $O(\log Z+\log(\ell+d+\log_2L))$.

Since the search starts at the least power of eight at least
$\max(\Ls,\ell^6)$ and multiplies by eight, it reaches every larger
power of eight unless it has stopped before. It therefore suffices to
show that both stopping tests hold at the following candidate.
Take $L$ to be the first power of eight at least $Z_L\ell^6$.
Then $Z_L\ell^6\le L<8Z_L\ell^6$. Both
$\lfloor\cs L^{1/3}\rfloor$ and $\lfloor\cs L\rfloor$ exceed
\eqref{c:eq:budgetcoarse}. To check this without an asymptotic assertion,
use $\log_2L\le3Z^8+6\ell+3$ and
$L^{1/15}\le2Z_L^{1/15}\ell^{2/5}$. The smallest power of $Z_L$
on the left is $Z_L^{1/3}$, and
\[
 Z_L^{1/3-1/15}=2^{(4/5)Z^8},\qquad
 Z_L^{1/3}=2^{Z^8}.
\]
These exceed $2^{1028}Z^6(3Z^8+20)$ for $Z\ge2^{2048}$;
repeated squaring, or the binomial lower bound for $2^N$, verifies
this last integer inequality. Also $\ell^2\ge\ell$ and
$\ell^2\ge\ell^{2/5}$. These comparisons prove both stopping tests
at this candidate. The search must terminate no later, so
\begin{equation}
 L\le8Z_L\ell^6.                                   \label{c:eq:Lexplicit}
\end{equation}
Inserting this bound in~\eqref{c:eq:dcoarse} gives
$d\le Z^3Z_L\ell$ after increasing only the displayed conservative
constant;~\eqref{c:eq:tcoarse} then gives $t_{\rm b}\le Z^8Z_L\ell$.
This proves~\eqref{c:eq:polylogscales}.

All searched integers have $O(\ell)$ bits, with fixed constants as
above. The tested powers $2^{\lfloor\cs \ell_{\rm buf}/L\rfloor}$,
$2^{\lfloor\cs L^{1/3}\rfloor}$ and $2^{\lfloor\cs L\rfloor}$ have
polynomially many bits in $\ell$ (at most $O(\ell^6)$, with a fixed
constant of order $Z_L$). The number of doublings is $O(\ell)$, also with such a
fixed constant. Ordinary exact rational arithmetic therefore has
polynomial-in-$\ell$ cost. In particular no failed candidate requires
a large ambient matrix or continuum eigenproblem to be constructed.
Finally $\alpha^{-1}=2^{1026}L^{-2}8^d$,
$\tau^{-1}<2^{25}L^4 3^{2d}$, and
$L_b=64mN_\Omega/\tau$ are polynomial in $m$ by the bounds just proved.
The same is true of every later precision budget.
\end{proof}

The finite scalar data and templates are calibrated only after this
parameter search. The actual material approximation is fixed before
sampling $\sigma$. Consequently the counting argument in
Lemma~\ref{c:lem:detuning} is applied to one fixed self-adjoint operator,
not to an operator chosen as a function of its own random detuning.

\section{Finite calibration without a continuum-solver oracle}
\label{c:sec:calibration}

This section gives algorithms for $z_h,\mu_h,u_h$ and the three
interactions. Their error estimates refer to the infinite operators
already defined, not to an unproved convergence property of a mesh.
Only the fixed corner matrix is inverted densely. Every other operation
is a finite Fourier operation, a convergent geometric series, or a
fixed-precision elementary function calculation.

\subsection{Elementary certified arithmetic and the edge energy}
Rational enclosures for $\pi$ follow from
$\pi=16\arctan(1/5)-4\arctan(1/239)$ and the alternating arctangent
series. The next term bounds its remainder. Square roots are enclosed
by integer bisection or integer square root. Exponentials and sine and
cosine use range reduction, Taylor series, and factorial geometric-tail
bounds. On arguments of magnitude $O(N_{\rm b})$ these procedures take a fixed
polynomial number of bit operations in $N_{\rm b}$ for $N_{\rm b}$ accurate bits.
Arguments outside this range are reduced, or the exponentially small
quantity is enclosed directly. The same statements for $\chi$ were
proved in Lemma~\ref{c:lem:gevrey}. Throughout, intervals are rounded
outward to a dyadic grid. Exact Gaussian elimination on a fixed-size rational matrix has
polynomial bit cost: clearing denominators uses only a fixed number
of entries, and every intermediate determinant has fixed order.
For a real symmetric matrix its ordered eigenvalues can be enclosed
without a root-separation assumption. At a rational trial value,
compute the inertia after subtraction of that value by symmetric
elimination. A nonzero diagonal gives a one-by-one pivot. If all
remaining diagonals vanish but an off-diagonal entry is nonzero,
its two-by-two block has one positive and one negative eigenvalue
and is an invertible pivot. An entirely zero remaining matrix
contributes only zeros. Congruence preserves the inertia at each
step. Counting positive eigenvalues therefore decides on which side
of the trial value the largest eigenvalue lies, including equality.
Bisection of an initial row-sum bound gives an interval of width
$2^{-N_{\rm b}}$ in $O(N_{\rm b})$ steps, with polynomial bit cost in the entry
lengths and $N_{\rm b}$. Interval matrix errors are added by the min--max
norm bound. These are explicit finite algorithms, not assumptions
about floating-point software.

For a local fiber, retain all 24 components of the eight central
aliases. Let $D_{\rm free}(q,\mathcal E)$ be the complementary free diagonal operator.
For $|q_j|\le3q_0$ and $\mathcal E$ in the edge window or within $2h^2$ of
$E_c$, its inverse is explicitly~\eqref{c:eq:freeinverse}, with norm
at most one. The complementary material inverse is
\begin{equation}
 \big(D_{\rm free}-\varrho V_{cc}\big)^{-1}
 =\sum_{a=0}^{\infty}
       (D_{\rm free}^{-1}\varrho V_{cc})^aD_{\rm free}^{-1}.              \label{c:eq:neumanncalibration}
\end{equation}
The tail after $K$ terms has norm at most $2(7\varrho)^K\le2^{1-K}$.
Each term changes a Fourier alias by at most $K$. The 24 by 24
central Schur matrix is consequently computable with a certified
operator-norm error at most $2^{10-K}$ using only aliases within
$K+2$ steps of the central set.

At $q=0$, the largest eigenvalue of that Schur matrix is zero exactly
at $E_c$ in the interval supplied by Lemma~\ref{c:lem:edge}. In that
interval all other Schur eigenvalues are negative. Its derivative
with respect to $\mathcal E$ is
\[
 -I-A_{{\rm per},0c}(A_{{\rm per},cc}-\mathcal E)^{-2}A_{{\rm per},c0},
\]
with the central and complementary alias blocks of the fiber as subscripts,
so its quadratic-form derivative lies between $-2I$ and $-I$.
The largest Schur eigenvalue therefore decreases at a rate between
one and two. This follows also by the min--max principle applied to
the two Schur matrices, and does not assume differentiability of an
ordered multiple eigenvalue.

Choose a rational $r_c$ with
$|r_c-E_0-\varrho\lambda_*|\le2^{15}\varrho^2$, using 256-bit
enclosures of $\pi$ and $\sqrt{19}$. The initial bracket
$[a,b]=[r_c-2^{17}\varrho^2,r_c+2^{17}\varrho^2]$ has width
$2^{-182}$ and strictly contains the band enclosure.
Put $w=b-a$ and $c=(a+b)/2$. Approximate the central Schur
matrix at $c$ to operator-norm error at most $w/64$, then test the
rational matrix for negative definiteness by exact elimination.
The slope bound gives $E_c\le c+w/64$ when the test is true and
$E_c\ge c-w/64$ otherwise. Thus use the overlapping brackets
\begin{equation}
 [a,b]\longmapsto
 \begin{cases}
 [a,c+w/64],&\text{negative definite},\\
 [c-w/64,b],&\text{otherwise}.
 \end{cases}                                             \label{c:eq:edge-bracket-step}
\end{equation}
Both branches have width $33w/64$, including when $c=E_c$.
Since $(33/64)^2<1/2$, $2N_{\rm b}$ steps give width at most
$2^{-N_{\rm b}}$. If $a=a_0/d_0$ and $b=b_0/d_0$, the first branch stores
$(64a_0,31a_0+33b_0,64d_0)$ and the second stores
$(33a_0+31b_0,64b_0,64d_0)$; every step adds at most six magnitude bits.
At step $k$, matrix error $2^{-(188+k)}$ is at most $w/64$, so the
working precision and endpoint lengths grow linearly with the number
of steps. Taking $N_{\rm b}$ a fixed amount larger than
$4\log_2L$ and rounding supplies~\eqref{c:eq:z} in polynomial bit time.

\subsection{Finite approximation of the local inverse}
On the real chart at $\mathcal E=z_h$, the exact full fiber inverse has norm
at most $4L^2$. Its central Schur inverse is a compression of the
full inverse and has the same bound. Approximate the Schur matrix
to norm error at most $2^{-N_{\rm b}}/(2^{20}L^4)$, invert the rational
approximation, and enclose its inverse using the inverse-perturbation
identity. The approximate inverse exists because its perturbation
relative to $4L^2$ is less than one half. Exact elimination or a
rational inverse followed by a residual test therefore succeeds.
The complementary inverse and the off-diagonal Schur factors use
\eqref{c:eq:neumanncalibration}. The block inverse formula then
approximates the full local inverse to norm error $2^{-N_{\rm b}}$ once
\begin{equation}
 K\ge N_{\rm b}+16\lceil\log_2L\rceil+2^{20}.               \label{c:eq:NeumannN}
\end{equation}
This deliberately weak bound covers all block products. Multiplication
by an enclosure of $\chi(q)$ gives the same precision for $R_0$.

The approximation is an operator-norm statement on the \emph{whole}
fiber, including arbitrarily high longitudinal aliases. No high
frequency has been removed on the assertion that its curl energy is
large. Formula~\eqref{c:eq:freeinverse} treats those directions exactly
before the small-potential series is truncated. On a finitely supported
Fourier input, the truncated formula produces finitely supported output;
its alias radius increases by at most $2K+4$.

All computations at a rational Bloch-grid index involve $\pi$ times
rational numbers and the already enclosed $z_h$. Noncentral free
denominators are bounded away from zero. The central inverse is
bounded by $4L^2$. Thus $N_{\rm b}+O(\log L+\log(2+|\nu|))$ working bits
suffice for the free entries, and a fixed multiple of that number
suffices for the fixed-size Schur calculation. An interval straddling
the edge of the cutoff support does not require a sign decision:
the inverse is valid on the larger chart $|q_j|\le3q_0$, while the
cutoff value is enclosed using its smooth defining formula.

\subsection{Periodization, with a norm error bound}
For calibration only, take a torus of side $2M_{\rm tor}$, with $M_{\rm tor}$ a multiple
of four. Place either one well, or the two wells needed for a
particular interaction, well inside it. Use a centered placement
so that the distance between these supports and any nontrivial
periodic image is at least $M_{\rm tor}/2$. A sufficient geometric condition
is $M_{\rm tor}\ge16(L+2Ld\max B_{jj}+1)$.

The periodized kernel is the sum of all translates by $2M_{\rm tor}\Z^3$.
Lemma~\ref{c:lem:cellbounds} makes that sum converge in the relevant
cell-block norms. On one or two well supports, the norm of the image
error is at most
\begin{equation}
 \Cs^2 L^{10}
       \exp\left(-\frac{\cs}{8}\sqrt{M_{\rm tor}/L}\right).
                                                               \label{c:eq:perioderror}
\end{equation}
To verify the factor, each support has at most $2^6L^3$ period cells.
For image index $j\ne0$ the cell distance is at least
$M_{\rm tor}|j|/4$. Sum the bound~\eqref{c:eq:cellgevrey} using shells with at
most $2^{10}(s+1)^2$ indices at $\ell^\infty$ distance $s$.
The elementary substitution $u=\sqrt s$ bounds the tail by a fixed
polynomial in $\cs^{-1}$ and $M_{\rm tor}/L$, times
$e^{-(\cs/8)\sqrt{M_{\rm tor}/L}}$. Weakening the exponent once absorbs that
polynomial; the remaining constants are less than $\Cs^2$.
The row and column Schur bounds give~\eqref{c:eq:perioderror}.

Choose the least multiple of four with the geometric condition and
\begin{equation}
 M_{\rm tor}\ge\Cs^4 L(N_{\rm b}+d+\lceil\log_2L\rceil+2)^2.
                                                               \label{c:eq:Mcalibration}
\end{equation}
Then~\eqref{c:eq:perioderror} is less than $2^{-N_{\rm b}-20}$.
The side $2M_{\rm tor}$ is an auxiliary calculation domain, \emph{not} a
physical internal boundary or an additional device volume. Its
computational cost is counted below. The periodic kernel is exactly
diagonalized by the finite set of Bloch momenta $q=\pi j/M_{\rm tor}$ modulo
$\pi\Z^3$, with the antiperiodic corner included because $M_{\rm tor}$ is even.
For each such momentum its alias space is still infinite until the
certified Neumann application above is made. This is a direct
Fourier-series identity for periodization, not a quadrature assumption.

\subsection{Multiplication and the isolated eigenvector}
On this torus approximate $q_w(hx)$ by a real trigonometric polynomial
$\widetilde q_w(r)$. The tail estimate~\eqref{c:eq:gevreytail} gives a uniform
error at most $2^{-N_{\rm b}}$ with Fourier index radius
\begin{equation}
 N_Q\le\Cs^4(1+M_{\rm tor}/L)(N_{\rm b}+\lceil\log_2M_{\rm tor}\rceil+2)^2.
                                                               \label{c:eq:NQ}
\end{equation}
For clarity, the Fourier coefficient of the scaled compact bump is
$(L/(2M_{\rm tor}))^3$ times its Fourier transform at $\pi L\nu/M_{\rm tor}$ with the
fixed normalization factor. Summing the tail outside $|\nu|_\infty>N_Q$
is bounded by the rescaled integral in~\eqref{c:eq:gevreytail}, plus
one layer of cubes; increasing the fixed factor in~\eqref{c:eq:NQ}
covers that layer. The coefficients are computed by a periodic
trapezoidal rule on a grid with a fixed larger index radius. Its
alias error is bounded by the same Fourier tail. There are
$O(N_Q^3)$ nodes. Each bump value has the certified evaluation
procedure in Lemma~\ref{c:lem:gevrey}. Rounding coefficients with an
extra $O(\log N_Q)$ guard bits controls their total uniform error.
Thus multiplication is approximated in \emph{operator norm}, rather
than just on a set of sampled fields.

The exact torus compression with the compact bump acts as zero
outside its support. On that support it differs from $S_h$ by
\eqref{c:eq:perioderror}. Replacing both bump factors by $\widetilde q_w$, and
using the norm-$2^{-N_{\rm b}}$ inverse approximation, changes its norm by
at most $2^{10}L^2 2^{-N_{\rm b}}$. All these comparisons take place in the
same $L^2$ space after extension by zero; the extra zero subspace
on the torus creates no ambiguity about the isolated largest
positive eigenvalue.

Here is an explicit polynomial-in-precision eigenvector algorithm.
Scale $S_h$ to $T_h=2^{90}L^{-2}S_h$. Lemma~\ref{c:lem:effectivemass}
gives
\begin{equation}
 2^{-38}\le\lambda_{\rm top}(T_h)\le1,
 \qquad
 \max_{\lambda\ne\lambda_{\rm top}}
       \frac{|\lambda|}{\lambda_{\rm top}}
       \le1-2^{-104}.                              \label{c:eq:powergap}
\end{equation}
For the negative spectrum use the separated-band bound
$2^{112}$ on $R_0$, while the selected part is nonnegative.
Since $L\ge\Ls$, its scaled negative norm is less than
$2^{-38}(1-2^{-104})$. The positive complementary spectrum uses
the gap $2^{-192}L^2$ and the upper eigenvalue bound.

Start with the normalized vector in Lemma~\ref{c:lem:effectivemass},
whose overlap is at least $2^{-48}$. Apply an approximation to
$T_h$ successively
\begin{equation}
 t_{\rm pw}=\left\lceil2^{106}(N_{\rm b}+200)\right\rceil            \label{c:eq:powersteps}
\end{equation}
times, then normalize. The exact normalized power has error at most
$2^{-N_{\rm b}-20}$ by~\eqref{c:eq:powergap}. Its unnormalized norm is at
least $2^{-48-38t_{\rm pw}}$. Request each operator application and the
starting vector to absolute accuracy
\begin{equation}
 2^{-N_{\rm work}},\qquad
 N_{\rm work}=N_{\rm b}+40t_{\rm pw}+\lceil\log_2(t_{\rm pw}+1)\rceil+256.
                                                               \label{c:eq:workbits}
\end{equation}
The telescoping product bound, using $\|T_h\|\le1$, makes the total
unnormalized error at most
$2(t_{\rm pw}+1)2^{-N_{\rm work}}$. Dividing by the proved lower norm and
using the elementary Lipschitz bound for normalization leaves less
than $2^{-N_{\rm b}-10}$. Torus, multiplier, and Neumann errors are allocated
inside the same per-application budget. Consequently the returned
normalized vector is within $2^{-N_{\rm b}}$ of $u_h$ in $L^2$, with its
sign fixed by the starting vector. Taking the real part and
renormalizing uses another two guard bits. A Rayleigh quotient,
computed by one more certified application, then encloses $\mu_h$
to absolute error $2^{-N_{\rm b}}$ after adding $O(\log L)$ bits.
This algorithm is intentionally slow in its \emph{fixed} gap constant;
its dependence on requested precision remains polynomial.

Every vector in this procedure is a finite Fourier polynomial. A
truncated inverse application increases the ordinary torus Fourier
radius by at most $(2K+4)M_{\rm tor}$, and multiplication by $\widetilde q_w$ increases it
by at most $N_Q$. The stored initial carrier has radius at most $3M_{\rm tor}$. For its initial bump
approximation, a uniform multiplier error $e$ produces an $L^2$
error at most $2^7(1+M_{\rm tor}/L)^2e$ after the $h^{3/2}$ scaling.
The exact unnormalized starting vector has norm at least $1/4$
by the overlap calculation in Lemma~\ref{c:lem:effectivemass}.
Thus using uniform error at most
$2^{-N_{\rm work}}/[2^{12}(1+M_{\rm tor}/L)^2]$, and allocating the
same accuracy to the fixed carrier coefficients, supplies the
requested normalized starting-vector precision. These are only
$O(\log M_{\rm tor})$ extra guard bits after $M_{\rm tor}$ is chosen. Each multiplier
and inverse application is likewise assigned its stated operator
error before summation; the explicit factors $4L^2$ and $2^{90}L^{-2}$
require only $O(\log L)$ further bits.
Thus after all applications the radius is bounded by
\begin{equation}
 N_{\rm all}\le3M_{\rm tor}+(t_{\rm pw}+2)(2N_Q+(2K+4)M_{\rm tor}).
                                                               \label{c:eq:allfreq}
\end{equation}
All $M_{\rm tor}^3$ Bloch-grid entries, all Fourier coefficients, and even naive
quadratic-time convolutions therefore cost a fixed polynomial in
$L,d,N_{\rm b},\log L$. With $L=O(\ell^6)$ and $N_{\rm b},d=O(\ell)$, this is a
fixed polynomial in $\ell$. There is no dense solve of dimension
comparable to $m$ in this calculation.

\subsection{Interaction and tuning calibration}
Use the two-well torus and the same mode approximation to evaluate
\[
 h_j(d)=\langle a_h,R_0\mathcal T_{2LdBe_j}a_h\rangle.
\]
Fourier translation is multiplication by explicitly computable unit
phases, and the inner product of two finite Fourier vectors is a
finite sum. The inverse norm $4L^2$ bounds the error from either
source-mode approximation;~\eqref{c:eq:perioderror} bounds images;
and~\eqref{c:eq:NeumannN} bounds the inverse error. Requesting each
of these errors at most $\zeta^2/(2^{12}L^2)$ suffices for the
absolute interaction error $\zeta^2/16$. There is no division by a
numerically guessed coupling. Its nonzero lower bound is already
proved in Lemma~\ref{c:lem:interactions}.

The same calculation supplies $\mu_h$ to $\zeta^2/16$ and the mode
or source mode to any later required $L^2$ tolerance. The targets
\eqref{c:eq:T}, their reciprocals, and the dyadic $v_i$ are then formed
using~\eqref{c:eq:vtuning}. The three direction calibrations and the
single-well data are shared by all sites. The targets $\Theta_{ii}$ take
only finitely many values, $\alpha(\Gamma_{ii}-\deg i)$ with
half-integer $\Gamma_{ii}\in[-3,3]$ and $\deg i\le3$ for terminals
and three values $-\beta_j^2/\alpha$ for mediators, so only $O(1)$
distinct $v_i$ are computed. Every requested bit count
is $O(\ell)$ with fixed constants, by Lemma~\ref{c:lem:parameters}.
This proves both termination and the cost of every input-dependent
calibration. In particular, the effective-mass estimate, which is only
$O(h)$, is \emph{not} used as an inverse-polynomial approximation to
the shrinking couplings. It establishes isolation and overlap; the
certified calculations in this section refine the actual auxiliary
operator to the much smaller required error.

\section{Finite material, source, and readout representations}
\label{c:sec:representation}

\subsection{The representation format}
The box sides, centers, and common radius $L$ are integers. In a ball
use the local coordinates $y=(r-r_i)/L$ and $t=|y|^2$. The material
is exactly~\eqref{c:eq:actualmaterial}. Its background is the fixed
finite expression in cosines in~\eqref{c:eq:background}; its only
input-dependent background coefficient is the dyadic $z_h+\sigma$.
A ball profile is specified by dyadic coefficients in the shifted
Chebyshev basis $T_j(2t-1)$ for $\widetilde p_{\rm r}(t)$, followed by the stated
square. This factored representation avoids large coefficients from
unnecessary conversion to monomials.

A source or readout template has the form
\begin{equation}
 \widehat a(r)=\widetilde p_{\rm r}(|y|^2)\widetilde P(y)\one_{|y|<1},
                                                               \label{c:eq:templateformat}
\end{equation}
where $\widetilde P$ has three components and is a tensor-product Chebyshev
polynomial on $[-1,1]^3$, with dyadic coefficients. A source is a
sum of translates of this template with dyadic amplitudes; a readout
uses the real part of its $L^2$ pairing with $E$. Values on the sphere have
no effect on these functions as $L^2$ objects. Both the formula and
its finite coefficients, rather than a quadrature table with an
unstated error, define the output function.

The format has a direct evaluation algorithm. Integer translations,
scaling, squared radius, polynomial recurrences, and the fixed cosines
are evaluated to a requested precision by outward arithmetic. On
their real domains, $|T_j|\le1$ and their derivatives are bounded
by $j^2$. Polynomial evaluation therefore needs only logarithmically
many guard bits in the number of coefficients and their magnitude.
The indicators specify exact semialgebraic balls. Away from their
boundaries membership is decided by the rational squared-radius
comparison; arbitrary-real point evaluation at a boundary is not
used as a computation primitive. The continuum instance is defined
by these exact functions independently of any point evaluator.

\subsection{A polynomial profile with endpoint matching}
Let $p_{\rm r}(t)=e^{-1/[2(1-t)]}$ for $0\le t<1$, with $p_{\rm r}(1)=0$.
The proof of Lemma~\ref{c:lem:gevrey}, now in one variable, bounds
its derivatives by $2^{20K}(K!)^2$. Composition with
$t=(1+\cos\varphi)/2$ gives a periodic even function with bounds
$2^{30K}(K!)^2$; this follows by the product rule, or inductively
by applying $-(\sin\varphi/2)\partial_t$ and differentiating its
bounded trigonometric coefficients. Integrating its Fourier
coefficients by parts $K$ times and choosing $K$ proportional to
$\sqrt j$ gives
\begin{equation}
 |c_j|\le 2^{100}\exp(-2^{-40}\sqrt j)
                                                               \label{c:eq:chebtail}
\end{equation}
for its Chebyshev coefficients. The same bound with a fixed
polynomial factor in $j$ controls the once differentiated series,
since $\|T_j'\|_\infty\le j^2$.

Truncate at degree $D_{\rm pr}\le\Cs(N_{\rm b}+1)^2$, choosing the least larger
integer for which the explicit tail bound is below $2^{-N_{\rm b}-20}$.
The sum of $j^2e^{-2^{-40}\sqrt j}$ above $D_{\rm pr}$ is bounded by its
first term plus the corresponding integral; the substitution
$u=\sqrt j$ makes this an elementary exponential-polynomial bound.
Compute the coefficients by the discrete cosine transform on a
fixed larger grid. The aliased-coefficient error is again bounded
by~\eqref{c:eq:chebtail}, and the samples have certified elementary
function evaluations. Round coefficients with
$N_{\rm b}+O(\log(D_{\rm pr}+2))$ fractional bits. If $\widetilde p^{\rm raw}_{\rm r}$ is the result,
replace it by
\begin{equation}
 \widetilde p_{\rm r}(t)=\widetilde p^{\rm raw}_{\rm r}(t)-\widetilde p^{\rm raw}_{\rm r}(1)
                     -(\widetilde p^{\rm raw}_{\rm r})'(1)(t-1).     \label{c:eq:endpointmatch}
\end{equation}
The endpoint values are exact dyadic sums because
$T_j(1)=1$ and $T_j'(1)=j^2$; for the shifted basis the
endpoint derivative is $2j^2$. The two correction coefficients have
$O(N_{\rm b}+\log D_{\rm pr})$ bits after the allocated rounding. The exact identities
$\widetilde p_{\rm r}(1)=\widetilde p_{\rm r}'(1)=0$ hold, and the uniform function and derivative
errors are at most $2^{-N_{\rm b}}$. In particular the zero extension of
$\widetilde p_{\rm r}(|y|^2)$ is $C^1$, while its square is nonnegative. All these
operations take polynomial-in-$N_{\rm b}$ bit time.

Since the well supports are disjoint, the material approximation
error is at most $3\max_i v_i\,2^{-N_{\rm b}}$, with no factor $J$.
Choose $N_{\rm b}$ so this is within~\eqref{c:eq:materialprecision}. This is
$O(\ell)$ bits. The ball-profile degree is therefore $O(\ell^2)$.

\subsection{Finite approximation of the calibrated mode}
The calibration algorithm returns a finite real Fourier polynomial
$\widetilde u(r)$ approximating $u_h$ on the ball in $L^2$, with the accuracy
proved in Section~\ref{c:sec:calibration}. On the local cube put
$\widetilde u_L(y)=\widetilde u(Ly)$. Each coordinate frequency has magnitude at most
$\pi N_{\rm all}L/M_{\rm tor}$. Its coefficient $\ell^1$ norm is at most the
square root of its number of Fourier coefficients times its $L^2$
norm, with the explicit torus normalization. These are polynomially
bounded in $L,d,N_{\rm b}$.

A finite trigonometric polynomial of maximum coordinate frequency $k_f$
and coefficient $\ell^1$ norm $a_f$ has a tensor Chebyshev approximation
of coordinate degree
\begin{equation}
 D_f\le 16(k_f+N_{\rm b}+\log_2(2+a_f)+1)                  \label{c:eq:chebdegree}
\end{equation}
with uniform error $2^{-N_{\rm b}}$, after a fixed increase in $N_{\rm b}$.
Here is a direct bound. Substitute $y_j=(s_j+s_j^{-1})/2$ and use
$|s_j|=2$. The imaginary part of $y_j$ has magnitude at most $3/4$,
so the function on this polyannulus is bounded by
$a_f e^{(9/4)k_f}$. Cauchy's Laurent coefficient estimate gives a
factor $2^{-j}$ in each coordinate. Summing the omitted geometric
tails proves~\eqref{c:eq:chebdegree}. A tensor discrete cosine transform
with a fixed larger degree gives coefficients with a certified
alias bound by the same estimate. True coefficients are at most
eight times the real-domain supremum, so their integer bits are
bounded by $O(\log(2+a_f))$, not by $k_f$. Rounding each coefficient
with $N_{\rm b}+O(\log(D_f+2))$ fractional bits controls the total uniform
error. This gives a polynomial-cost construction, not only an
existence result for a polynomial approximation.

Use this procedure with a uniform tolerance smaller than the desired
local $L^2$ tolerance divided by $8L^{3/2}$. Multiply the resulting
$\widetilde P(y)$ by $\widetilde p_{\rm r}(|y|^2)$ as in~\eqref{c:eq:templateformat}.
The multiplier approximation is in uniform norm, the calibrated
mode approximation is in $L^2$, and the cube has volume $8L^3$.
Their triangle inequality therefore gives a certified bound
\begin{equation}
 \|\widehat a-a_h\|_{L^2(\mathbb B_L)}\le e_a          \label{c:eq:atesterror}
\end{equation}
for any prescribed inverse-polynomial $e_a$.
This does not require unproved regularity or spectral convergence
of a piecewise-polynomial Maxwell discretization.

For the chosen parameters, the index bounds are
$M_{\rm tor}=O(\ell^8)$, $N_Q=O(\ell^4)$, $K=O(\ell)$,
$t_{\rm pw}=O(\ell)$, with the conservative stored-radius bound
$N_{\rm all}=O(\ell^{11})$. For the physical frequency, use
$LN_Q/M_{\rm tor}=O(\ell^2)$ directly in~\eqref{c:eq:allfreq}:
$LN_{\rm all}/M_{\rm tor}=O(\ell^8)$, so $k_f=O(\ell^8)$.
The degree in~\eqref{c:eq:chebdegree} is therefore
$O(\ell^8)$, giving $O(\ell^{24})$ tensor coefficients and
$O(\ell^{25})$ encoded bits per template. Each coefficient, including
its integer part, has $O(\ell)$ bits. The intermediate
exact arithmetic or guard precision may have a fixed polynomial
number of bits in $\ell$, which is counted in the time bound.

\subsection{Source rounding and stable observables}
Define the rational budget
\begin{equation}
 e_f=\min\{e_0/4,\ 2^{-260}\alpha L^{-2}\},
 \qquad
 e_a\le\min\{e_f/4,e_0/4\}.                 \label{c:eq:sourcebudgets}
\end{equation}
Construct one shared template satisfying~\eqref{c:eq:atesterror} and
$\|\widehat a\|\le2$. Enclose $\alpha/\mu_h$ to absolute error
$e_f/16$. Multiply its dyadic approximation by the short rational
nonzero entries of $g_\Theta$ and round each product to a dyadic amplitude
with error at most $e_f/(16m)$. There are at most $m$ such amplitudes.
Call the resulting finite source $f$. Disjoint supports and
$\alpha/\mu_h\le1$ give
\begin{equation}
 \|f-f_*\|\le e_f \norm g\le e_0\norm g/4,
 \qquad 0<\|f\|\le2\norm g.                              \label{c:eq:finite-source-error}
\end{equation}
For example the template error contributes at most $e_f\norm g/4$, the
shared scalar error at most $e_f\norm g/8$ using $\|\widehat a\|\le2$,
and the amplitude rounding at most
$2\sqrt m\,e_f/(16m)\le e_f\norm g/8$.
The lower source bound in~\eqref{c:eq:auxnorms} and the second term
in $e_f$ ensure that the total error is at most one sixteenth of
$\|f_*\|$, so $f$ is genuinely nonzero. It is a finite sum of
bounded compactly supported functions, hence belongs to $L^2$.
No derivative of an indicator is part of its definition.

The ideal observable at well $i$ is
\begin{equation}
 z_i(E)=v_i\langle E,a_i\rangle_{L^2}.
                                                               \label{c:eq:idealreadout}
\end{equation}
It is a bounded $L^2$ functional of norm at most one. The finite
observable replaces $a_i$ by the translated template $\widehat a_i$.
Its norm is at most two. Because supports are disjoint, the norm
of the \emph{entire} template-error map is at most $e_a$:
\[
 \sum_i|\langle E,\widehat a_i-a_i\rangle|^2
 \le e_a^2\sum_i\|E\|_{L^2(\mathbb B_L+r_i)}^2
 \le e_a^2\|E\|^2.
\]
This is the stability of the prescribed volume readout. There is
no point evaluation of an $H(\curl)$ field.

Request the real part of each raw finite observable to absolute
error at most
\begin{equation}
 h_y=\frac{\ee}{2^{24}J_\Gamma\max(1,J)}.            \label{c:eq:momentaccuracy}
\end{equation}
Multiply by the computed $v_i$ and round with an additional absolute
error at most $h_y$. A dyadic grid no coarser than $h_y$ fixes the
message format. Reserve a signed integer range larger than $8L_bm$.
The decoder consumes only this fixed number of bits per message,
rejecting a malformed message or overflow. The integer and fractional
bit counts are $O(\ell)$ on every draw.

Producing the moments to this accuracy is part of the continuum
solver's output obligation, and its cost is included in the solver
bound below. No finite-time integration oracle for an arbitrary
unrepresented $L^2$ function is granted to the decoder. The theorem
specifies a finite list of stable observables, explicit kernels, and
finite accuracies; the decoder reads and processes exactly those bits.

\subsection{Geometric and profile resolution}
The only material-support interfaces are the explicit spheres and the
external box. Their radii are $L\ge1$, their mutual gaps are at least
$2Ld\min_j B_{jj}-2L\ge L$, and their boundary buffer is at least
$\ell_{\rm buf}\ge L$. The fixed microscopic background has period two. There
are no shrinking geometric supports, gaps, filaments, or layers.
The smooth within-ball variation is specified by the counted degree
and coefficients, rather than by an uncharged fine mesh.

For a quantitative resolution bound, a tensor polynomial of degree $D_f$
and coefficient sum $a_f$ has first derivative bounded by $3a_fD_f^2/L$
in physical coordinates on the local cube. The radial profile obeys
the analogous $4a_fD_f^2/L$ bound; the squared profile only changes the
fixed factor. Here $D_f$ and the coefficient sums are fixed polynomials
in $\ell$ (the radial profile is uniformly bounded with its first
derivative by the preceding construction). The source and readout
frequency and degree bounds above are also fixed polynomials in
$\ell$. Thus even an explicit representation refined at the associated
inverse-polylogarithmic resolution has only a polylogarithmic extra
cost per unit volume. The proof does not rely on such a refinement
for PDE accuracy. It relies on uniform multiplier error for material
and $L^2$ error for source and readout. Small-amplitude polynomial
oscillations or close level sets therefore hide neither a geometric
list nor an uncounted discretization requirement.

\section{Every-field recovery for the actual continuum operator}
\label{c:sec:recovery}

Fix all parameters, finite profiles, and source and test coefficients,
then sample the detuning according to Lemma~\ref{c:lem:detuning} for
$\overline\varepsilon_0$ in~\eqref{c:eq:actualmaterial}. Set
\begin{equation}
 A=\Cc-\overline\varepsilon,
 \qquad \delta=e_0=\frac{\ee}{2^{24}J_\Gamma L_b}.
                                                               \label{c:eq:finaldelta}
\end{equation}
The auxiliary estimates hold for every draw $0\le\sigma<\tau$.
Tests~\eqref{c:eq:buffer-test}--\eqref{c:eq:leak-test}, the material
budget~\eqref{c:eq:materialprecision}, and the finite source budget
prove
\begin{equation}
 \|AE_{\rm ref}-f\|\le e_0\norm g.                          \label{c:eq:comparison-residual}
\end{equation}
On the common event of Lemma~\ref{c:lem:detuning}, take any
$E\in X$ satisfying~\eqref{c:eq:physicalres}. By~\eqref{c:eq:finite-source-error},
\begin{align}
 \|E-E_{\rm ref}\|&\le L_b(\delta\|f\|+e_0\norm g)
                 \le3L_be_0\norm g,                     \label{c:eq:fieldcomparison}\\
 \|E\|&\le L_b(1+\delta)\|f\|\le3L_b\norm g.             \label{c:eq:fieldrange}
\end{align}
Here $\delta<1/2$. Both inequalities apply to every residual field,
not only the exact solution and not only a Bloch-localized field.
All polarizations and all unwanted finite-box modes are controlled
by the full inverse bound in~\eqref{c:eq:fieldcomparison}.

The ideal moments of $E_{\rm ref}$ equal those of $E_{\rm a}$ because
the cutoff is one on all well supports. Their terminal vector $z_{\rm t}$
satisfies~\eqref{c:eq:auxlogical}. At well $i$ the first-moment error is
$v_i\ip{E}{\widehat a_i}-v_i\ip{E_{\rm ref}}{a_i}
=v_i\ip{E-E_{\rm ref}}{a_i}+v_i\ip{E}{\widehat a_i-a_i}$, plus at most
$2h_y$ from rounding. The first term uses $\norm{a_i}\le1$ and
\eqref{c:eq:fieldcomparison}, the second the template-error map of
Section~\ref{c:sec:representation} with~\eqref{c:eq:atesterror}
and~\eqref{c:eq:fieldrange}; with $0<v_i<1$ and disjoint supports this
shows that the finite moments of $E$ differ from this vector by at most
\[
 3L_be_0\norm g+3L_be_a\norm g+2\sqrt J\,h_y
 \le\frac{6\ee\norm g}{2^{24}J_\Gamma}.
\]
Consequently, with $\widehat z$ the decoded moment vector,
\begin{equation}
 \|\Gamma\widehat z_{\rm t}-g\|
 \le\ee\norm g/1024+J_\Gamma\frac{6\ee\norm g}{2^{24}J_\Gamma}
 <\ee\norm g/2.                                         \label{c:eq:final-logical}
\end{equation}
Lemma~\ref{c:lem:logical} gives the desired rational answer to
$Mx=b$. Since the material, source, and test fields are real,
replacing a complex field by its real part cannot increase its
residual; requesting real parts of the moments implements exactly
that replacement. The argument also allows every set of moment
errors within the specified absolute bounds.

By~\eqref{c:eq:fieldrange} and $\|\widehat a_i\|\le2$, the raw moments
have magnitude at most $6L_b\sqrt m<8L_bm$. Thus the fixed message
range is adequate on the same event. No union over solutions or
outputs is needed. On bad draws the finite format and all costs
are unchanged; only correctness and the inverse bound are not
asserted.

\section{\texorpdfstring{Complete resource accounting and proof of Theorem~\ref{c:thm:main}}{Complete resource accounting and proof of the continuum realization theorem}}
\label{c:sec:cost}

\subsection{Physical volume and finite data}
The logical side lengths are $O(\sqrt{m\ell})$. Doubling the lattice,
inserting mediators, multiplying by $2LdB$, and adding the buffer give
\begin{equation}
 \max_j N_j=O\big(Ld\sqrt{m\ell}+\ell_{\rm buf}+L+1\big),
 \qquad
 |\Omega|=O(m^{3/2}\ell^{24}).                       \label{c:eq:volume-final}
\end{equation}
Here Lemma~\ref{c:lem:parameters} gives $L=O(\ell^6)$ and $d=O(\ell)$.
The sharper exponent from these bounds is $45/2$; 24 is a convenient
integer upper bound. The estimate counts the entire rectangular box,
including empty space between all routes and wells and the external
buffer. In vacuum wavelength units the volume is $|\Omega|/(2\pi)^3$;
using the shortest material wavelength changes this only by a fixed
factor because $1\le\varepsilon\le25$.

There are $J=O((m\ell)^{3/2})$ well records. Each records an integer
center, the common integer radius, a dyadic increment, and a pointer
to a counted polynomial profile. The background has a fixed finite
formula. Even if the common source and test templates are copied at
every site rather than shared, the degree bounds in
Section~\ref{c:sec:representation} give at most
$O(m^{3/2}\ell^{64})$ coefficient and index slots.
Only at most $m$ source amplitudes are nonzero. The readout has at
most $J$ messages, each of $O(\ell)$ bits, and the final rational
answer has $n\le m$ coordinates of $O(\ell)$ bits.

An explicit ambient format can enumerate all unit cubes of the box.
For each cube, the candidate nearby placement site is located by
integer division in the diagonal lattice $2LdB\Z^3$, followed by a
constant number of neighbor checks and a lookup in the sorted occupied
site table. The balls are separated by much more than their radii,
so at most one can meet a given cube. A partial intersection is
represented by its exact quadratic inequality, not by a polygonal or
voxel approximation to the sphere. This enumeration costs
$O(|\Omega|\operatorname{polylog}m)$ and does not hide the background
volume in a sparse list.

\subsection{Numerical scales and operation counts}
The explicit inverse quantities are
\begin{equation}
 \alpha^{-1}=2^{1026}L^{-2}8^d,
 \quad \tau^{-1}<2^{25}L^4 3^{2d},
 \quad L_b=64mN_\Omega/\tau,
 \quad \delta^{-1}=2^{24}J_\Gamma L_b\ee^{-1}.
                                                               \label{c:eq:inverse-scales}
\end{equation}
Equations~\eqref{c:eq:logicalparameters} and
\eqref{c:eq:polylogscales} show that all are polynomial in $m$.
So are the reciprocals of every material, source, mode, interaction,
and message tolerance. All output coefficients and all integer
indices therefore require $O(\ell)$ bits.

The local-calibration bounds can be checked directly from the formulas:
$M_{\rm tor}=O(\ell^8)$, $N_Q=O(\ell^4)$, $K=O(\ell)$,
$t_{\rm pw}=O(\ell)$, and $N_{\rm all}=O(\ell^{11})$.
A coefficient array therefore has $O(\ell^{33})$ entries. Using
naive convolution takes $O(\ell^{66})$ scalar operations per product;
there are polynomially many such products and Bloch-grid operations.
The tensor polynomial conversion has $O(\ell^{24})$ coefficient slots. It may
be implemented by direct discrete cosine transforms and direct Fourier
sample evaluation, whose quadratic product counts are still fixed
polynomials in $\ell$. Dyadic interval operations, elementary series,
and the 24 by 24 rational inverses use polynomial bit time in the
working precision. The bit cost also includes allocation, address
arithmetic, coefficient normalization, and serialization of the
computed arrays. With schoolbook arithmetic, the resulting bound for
the shared calibration and template construction is
$O(\ell^{4096})$. Replication over the $O(m^{3/2}\ell^{3/2})$ sites
copies the already computed coefficient lists, at
$O(m^{3/2}\ell^{27})$ bit cost; it does not rerun calibration at
every site. These costs, the parameter searches, and the ambient
enumeration fit within~\eqref{c:eq:resources}.

For fully explicit constants, use $Z,Z_L$ from~\eqref{c:eq:ZAL} and set
\begin{equation}
 C_{\rm res}=2^{Z^{16}},\qquad c_A=c_\delta=C_{\rm res}.                         \label{c:eq:theoremconstants}
\end{equation}
Here is a conservative bookkeeping check. The fixed analytic constants
are at most powers $\Cs^{1024}$. The parameter bounds are
$L\le8Z_L\ell^6$, $d\le Z^3Z_L\ell$, and
$\ell_{\rm buf}/L\le Z^8Z_L\ell$. Allowing degree 4096 in those constants,
the bit target, and the local array sizes, the logarithm of their fixed
multiplicative constant is less than
$2^{20}Z^9<Z^{16}$. The same bound covers the fixed-startup arithmetic
needed to form the constants themselves. Thus $C_{\rm res}$ bounds the
multiplicative constants in~\eqref{c:eq:resources}.
The powers of $m$ in~\eqref{c:eq:inverse-scales} are at most
$2^{20}Z^4Z_L+\lceil c_\kappa\rceil+\lceil c_\eta\rceil+100$.
Their fixed coefficients and logarithmic factors can be absorbed in
$m^{C_{\rm res}}$ for $m\ge2$: use $\ell\le3m$, and absorb a constant factor $c_{\rm f}$ by
an extra exponent $\lceil\log_2c_{\rm f}\rceil$. The resulting exponent is
strictly smaller than $C_{\rm res}$. This proves the explicit bounds
$702L_b\le m^{c_A}$ and $\delta\ge m^{-c_\delta}$, not just an unspecified
polynomial assertion.

\subsection{The uniform encoder and decoder}
The complete algorithm, in its dependency order, is as follows.
\begin{enumerate}
\item Normalize and round the input; build the finite circuit, sign
mediators, and balanced three-dimensional route of
Sections~\ref{d:app:logical}--\ref{d:app:routing}. Compute $\Gamma,g,\overline K,\ee$ and the doubled
terminal--mediator sites.
\item Use Section~\ref{c:sec:parameters} to choose $L,d,\ell_{\rm buf}$, the actual
integer box, $\tau,N_\Omega,L_b$, and all error budgets. This step
writes only counts and dimensions until its stopping tests pass.
\item Verify the fixed corner certificate. Compute $z_h$, the isolated
single-well data, and the three actual auxiliary interactions to the
prescribed tolerances by Section~\ref{c:sec:calibration}. Compute and
round every $v_i$ by~\eqref{c:eq:vtuning}.
\item Form the finite radial profile and the finite source and readout
templates, with the global budgets in
Section~\ref{c:sec:representation}. These determine the actual
$\overline\varepsilon_0$ before detuning.
\item Draw the fixed number of fair bits for the midpoint grid in
Lemma~\ref{c:lem:detuning}, form $\sigma$, and output the actual material,
source, geometry, $\delta$, test kernels, message accuracies, and
message ranges. Enumerate the ambient unit cubes as just described.
\item The decoder consumes the fixed-width moment messages, forms
$\widehat z$, and performs the coordinate selection, rational rounding,
and rescaling in Section~\ref{d:app:logical-residual}. It performs no solve
of the original system or of the assembled auxiliary system. Its side
information is $O(m)$ words passed by the encoder: the coordinate-selection
map and the scale $\nu_M/(ma_M)$ of Section~\ref{d:app:logical-residual} (and, in
the Real-RAM bridge of Section~\ref{d:app:real-transfer}, the scales
$ma_F,\nu_F$).
\end{enumerate}
The decoder can scan a complete ambient-grid message stream, including
zero entries away from the prescribed readout sites. Its description
and this stream together use $O(m^{3/2}\ell^{25})$ bits. Parsing,
addressed selection and rational arithmetic on $O(\ell)$-bit words
give the separate bound
\begin{equation}
 T_{\rm dec}=O(m^{3/2}\ell^{27}),                         \label{c:eq:decoder-cost}
\end{equation}
which is included in the joint encoding and decoding allowance.
All loops and all refinement rules above have proved stopping bounds.
The physical medium and its fixed data are given explicitly; there is
no search for an unspecified photonic crystal or an assumed numerical
spectral certificate. The fixed verification and any enormous fixed
startup cost are included in~\eqref{c:eq:theoremconstants}.

If $b=0$, return $x=0$ directly. A required physical output can be
the unit cube with $\varepsilon=1$, $f=0$, and $\delta=1$, with a
decoder that returns zero. Its curl eigenvalues are zero or at least
$2\pi^2$, so $\spec(\Cc-1)\subset\{-1\}\cup[2\pi^2-1,\infty)$, and
\eqref{c:eq:maininverse} holds with $L_b=1$. The finitely many bounded input sizes
can similarly be solved by counted exact rational elimination and
accompanied by that dummy instance. In particular the tolerance bound
also holds at $m=1$. Their lengths and costs are bounded by the
fixed input constants; the inverse-event statement of Theorem~\ref{c:thm:main}
is stated for $m\ge2$.

\begin{proof}[Proof of Theorem~\ref{c:thm:main}]
Lemma~\ref{c:lem:logical} applies to every sparse input
satisfying~\eqref{c:eq:input}, without
any geometric restriction on its graph. Lemmas~\ref{c:lem:edge}--\ref{c:lem:numerator}
prove the continuum local spectral platform, including all fiber
polarizations and aliases, a simple tunable auxiliary mode, and a
nonzero source numerator. Lemma~\ref{c:lem:interactions} supplies the
three useful couplings and the full-block upper bound with a strict
rate margin. Lemma~\ref{c:lem:assembly} includes the entire infinite
defect complement and the sum of unwanted interactions.

The auxiliary channel is not substituted for the physical equation:
Lemma~\ref{c:lem:leakage} proves its actual leakage residual, and
Lemma~\ref{c:lem:cutresidual} puts it inside the finite PEC box in the
correct operator domain. The terminating choices in
Lemma~\ref{c:lem:parameters} make both errors small relative to the
shrinking useful interactions and to the full inverse bound. Sections
\ref{c:sec:calibration}--\ref{c:sec:representation} replace every input-
dependent datum by the stated finite representation with proved
error and cost.

Apply Lemma~\ref{c:lem:detuning} to that actual material. On its one
common event, the full continuum operator has~\eqref{c:eq:maininverse},
and Section~\ref{c:sec:recovery} proves successful decoding for every
field and every finite message error in the theorem. The bounds in
this section prove all physical, computational, and numerical resource
claims, on unsuccessful draws as well. This completes the proof.
\end{proof}

\section{Explicit analytic majorants for the continuum constants}
\label{c:app:ledger}

This section supplies the numerical envelope for the analytic constants
of Part~II. It replaces ``take a sufficiently large platform constant''
by finite arithmetic bounds.

\subsection{Fixed-order derivative envelope}
Set
\begin{equation}
 M_*=2^{16384},\qquad
 M_j=(j!)^4M_*^{(j+1)^2}\quad(0\le j\le32).
                                                               \label{c:eq:Mj}
\end{equation}
All derivatives below are directional derivatives along unit vectors;
mixed derivatives can be bounded by expansion in the three coordinate
directions, adding at most $3^j$, already absorbed in~\eqref{c:eq:Mj}.
The following list gives the underlying numbers before this envelope
is used.
\begin{center}
\begin{tabular}{@{}p{0.49\textwidth}p{0.42\textwidth}@{}}
\toprule
Quantity & Explicit bound used \\
\midrule
Local analytic momentum radius & $R_q=2^{-112}$ \\
Real cutoff chart scale & $q_0=2^{-140}$ \\
Hessian and quartic constant & $2^{99}I\le H\le2^{108}I$,
 $M_4=2^{350}$ \\
Nonselected fiber inverse on the small chart & $2^{112}$ \\
Cell eigenfunction Fourier $\ell^1$ norm & $32$ \\
Scaled denominator on a complex radius-two neighborhood & $2^{256}$ \\
Scaled source numerator on that neighborhood & $2^{128}$ \\
Pole contraction radius in transverse variables & $2^{-1024}$ \\
Final pole neighborhood & $|k_\perp|\le2^{-4096}$ \\
Regular-region denominator lower factor & $M_*^{-1}(1+|k|^2)$ \\
Fixed cutoff derivative bases & at most $2^{4200}$ \\
\bottomrule
\end{tabular}
\end{center}

The first five lines follow directly from Lemma~\ref{c:lem:edge} and
its Neumann construction. For the scaled denominator, insert the
Taylor series in $h^{-2}(\lambda(hk)-z_h)$.
On a complex polydisc of radius two around the shifted origin,
$|k|<4$, its quadratic term is at most $2^{112}$ and its quartic
remainder at most $2^{358}h^2<1$. The shift $|b_{j,h}|<1$ and the
constant term are covered by $2^{256}$. The same argument on a
slightly larger polydisc bounds derivatives by Cauchy's formula;
for instance $j!2^{4j+256}$ suffices up to order 32.

For the numerator, compact support gives
\[
 h^{3/2}\int|a_h(r)|\dd r\le\sqrt{|\mathbb B_1|}<3.
\]
On the stated complex neighborhood $|e^{\pm ihk\cdot r}|\le e^8$,
and the cell Fourier $\ell^1$ norm is at most 32. Thus each scaled
source factor has magnitude less than $2^{20}$, and their product
is below $2^{40}$; the table uses a larger number. Derivatives with
respect to scaled momentum introduce bounded $hx$ factors and
$h$ times the fixed cell derivative constants. Cauchy's estimates
therefore give the claimed derivative envelope, also on the whole
real scaled chart with the allowed imaginary shifts.

\subsection{The pole chart and its remainders}
Write the shifted denominator as $D(k_\parallel,k_\perp)$, with
$D_\parallel(0)=ic$, $c\ge1$, $D_\perp(0)=0$. Its second transverse derivative
is $Q_{j,h}$, where~\eqref{c:eq:polebounds} gives
$2^{98}I<Q_{j,h}<2^{110}I$. The preceding Cauchy bound implies, for
$|k_\parallel|,|k_\perp|\le1/8$,
\[
 |D(k_\parallel,k_\perp)-icu-\tfrac12k_\perp^TQ_{j,h}k_\perp|
 \le2^{280}(|k_\parallel|^2+|k_\parallel||k_\perp|+|k_\perp|^3).
\]
Consider the map
$k_\parallel\mapsto-[D(k_\parallel,k_\perp)-icu]/(ic)$. On
$|k_\perp|\le2^{-1024}$ and $|k_\parallel|\le2^{300}|k_\perp|^2$ it maps the latter
disc strictly into itself and has derivative in $k_\parallel$ of magnitude
less than one half. Indeed the derivative remainder is at most
$2^{282}(|k_\parallel|+|k_\perp|)$, and the quadratic main term is at most
$2^{110}|k_\perp|^2$. Iteration from zero converges geometrically and
uniformly on smaller complex polydiscs. Its limit $p(k_\perp)$ is
holomorphic, is the unique zero in that disc, and obeys the same
size bound. This is a direct quantitative proof of the implicit
function used in~\eqref{c:eq:pole}.

Cauchy's estimate on the transverse radius $2^{-1024}$ bounds the
third-order remainder in $p$ by $2^{1400}|k_\perp|^3$ on the smaller
half-radius polydisc. Differentiating $D(p(k_\perp),k_\perp)=0$ at zero gives
$p(k_\perp)=iw^TQ_{j,h}k_\perp/(2c)+O(|k_\perp|^3)$. The imaginary part of the quadratic
form is at least $2^{-15}|k_\perp|^2$. On $|k_\perp|\le2^{-4096}$ the cubic
bound is less than $2^{-2696}|k_\perp|^2$, giving the weaker stated
$\operatorname{Im}p(k_\perp)\ge2^{-256}|k_\perp|^2$.
The factor
\[
 d_1(k_\parallel,k_\perp)=\int_0^1D_\parallel(p(k_\perp)+t(k_\parallel-p(k_\perp)),k_\perp)\dd t
\]
is analytic and differs from $ic$ by less than $c/2$ on this
neighborhood. Cauchy division of the numerator can consequently
be performed by subtracting its value at $k_\parallel=p(k_\perp)$ and dividing
by $k_\parallel-p(k_\perp)$; the analytic remainder is the integral of the
corresponding derivative along that segment. These formulas bound
all derivatives without division by a small, unlocated zero.

One explicit choice of the auxiliary cutoffs is
$\psi(k_\parallel/2^{-2048})\prod_{j=1}^2\psi(k_{\perp,j}/2^{-4100})$,
with smaller copies for the core where a cutoff equals one.
All poles over the transverse support have magnitude less than
$2^{-7890}$, so they lie strictly inside the longitudinal core.
The derivative bases of these cutoffs are at most $2^{4200}$.
On the regular region outside the smaller pole neighborhood,
\eqref{c:eq:tube} at its limiting value gives a denominator lower
bound $M_*^{-1}(1+|k|^2)$. For small $k$ outside that neighborhood
use $|k|\ge2^{-4104}$ and the positive quadratic real part; for
$|k|\ge1$ use the symbol's quadratic lower bound. The interval
between these regions is covered by the same inequality after
weakening to $M_*^{-1}$.

The scaled symbol derivatives obey
\[
 |\partial^jD_h(k+ib_{j,h}e_j)|
 \le j!M_*(1+|k|)^{2-j}\quad(0\le j\le32)
\]
on that regular region. For $j\le2$ this is the quadratic Taylor
estimate. For $j\ge3$, differentiation contributes $h^{j-2}$;
on the scaled chart $h|k|\le4q_0$, so the Cauchy bound
$3\varrho j!R_q^{-j}$ gives the displayed inequality. For bounded
$k$ use the already proved radius-two bounds.
The recurrence
\[
 \partial^K(D^{-1})=-D^{-1}\sum_{a=1}^K
       \binom{K}{a}(\partial^aD)(\partial^{K-a}D^{-1})
\]
then bounds its derivatives by $(K!)^4M_*^{(K+1)^2}$ times
$(1+|k|)^{-2-K}$. The same recurrence applies to an invertible
matrix or bounded operator, preserving the order of the factors;
Cauchy's estimate is an alternative for the uniformly separated
complement. Leibniz's rule and the geometric-cutoff bounds give
$M_{32}^{16}$ as a common bound for the coefficient of
$(1+|k|)^{-2}$ in the sixteenth derivative of every regular
selected-band integrand.

For the other bands a useful explicit source bound avoids assuming
that their source coefficients have rapid decay. The rescaled Bloch
transform of $a_h$ has fiber norm at most a constant: Cauchy--Schwarz
over its at most $2^6L^3$ support cells, multiplied by $h^{3/2}$,
gives this bound. Its $j$th scaled-momentum derivative inserts
$h$ times a cell coordinate, of magnitude at most a fixed constant
on the support. Thus the same bound holds for each $j\le32$ with
the factor $M_j$. The nonselected resolvent and its derivatives
are uniformly bounded, and their scaled contribution has the
extra factor $h^2$. Integrating over a chart of volume $O(h^{-3})$
gives at most $M_{32}^{16}/h$. For the selected regular part,
$\int_{|k|\le4q_0/h}(1+|k|)^{-2}\dd k\le M_*/h$ gives the
same bound. This verifies the $L d^{-16}$ term in
\eqref{c:eq:greenremainder} for the entire complement.

For the transverse Gaussian remainder, the amplitude error is
bounded by $M_3^4|k_\perp|$ and the phase error by $M_3^4T|k_\perp|^3$.
Both exact and quadratic phases have imaginary part at least
$2^{-256}|k_\perp|^2$. The identity
$e^a-e^b=(a-b)\int_0^1e^{(1-t)b+ta}\dd t$ therefore preserves
the Gaussian majorant on the segment between them. Polar
integration gives
\[
 \int_{\R^2}|k_\perp|e^{-cT|k_\perp|^2}\dd k_\perp=O(c^{-3/2}T^{-3/2}),
\]
\[
 T\int_{\R^2}|k_\perp|^3e^{-cT|k_\perp|^2}\dd k_\perp
       =O(c^{-5/2}T^{-3/2}),\qquad c\ge2^{-257}.
\]
The gamma factors at these integer and half-integer arguments are
bounded by $2^{20}$ using their elementary factorial formulas.
The omitted Gaussian tail outside the fixed cutoff has a bound
of the same form, since $\sup_{T\ge1}\sqrt T e^{-aT}\le a^{-1/2}$
and its $a$ is explicitly at least $2^{-8500}$.
The full-line longitudinal tail differentiated 16 times is bounded
by $16!$ times a power of the reciprocal longitudinal-core radius,
at most $2^{40000}$. These bounds are all below $M_{32}^{32}$.
They justify every fixed remainder coefficient in
\eqref{c:eq:greenremainder}.

\subsection{Arbitrary-order cutoff and spatial estimates}
The preceding fixed-order envelope is not used as a claim about
arbitrarily high derivatives. Those estimates have their own explicit
factorial bounds. Recall $2^{16}/q_0=2^{156}$. In
\eqref{c:eq:almostanalytic}, $2^{156}|\beta|K\le1/4$. After using
$K!\le K^K$, the defect from holomorphy is at most
\[
 2^{45}2^{156}|\beta|(K+1)^2 4^{-K}
 \le2^{60}2^{-K}.
\]
Since $K\ge2^{-159}L$ for the nonzero shifts used here, weakening
the exponent gives $\Cs e^{-\cs L}$. For a fixed additional
$q$-derivative of order $j\le32$, each term of the truncated
Taylor sum is bounded by $4^{-a}$ times a polynomial of degree
$2j$ in its index $a$, multiplied by
$2^{45}2^{156j}(j!)^2$. Summing that geometric-polynomial series is
bounded by $(2j+2)!2^{10j+100}2^{156j}$, below $M_j^4$.
This verifies the uniform fixed derivative bounds for the shifted
cutoff, independently of the truncation order or of $L$.

For the summable cell estimate, the real local inverse has an
analytic strip of width at least $2^{-256}h$. To check it, split
off the rank-one pole and use the Taylor estimate and positive
quadratic form: a complex displacement with that width changes
$h^2+q^THq$ by at most one quarter of its magnitude, by
$2ab\le a^2+b^2$ and $H\le2^{108}I$. The separated spectrum
stays separated by the complementary Neumann argument.
Cauchy's formula, one or two local inverse factors, the Gevrey
cutoff, and bounded cell gauge derivatives give
\[
 \|\partial^K(\hbox{cell integrand})\|
 \le\Cs L^4(2^{512}L)^K(K!)^2\qquad(K\ge0).
\]
Integrating by parts in its cell Fourier coefficient and choosing
$K$ proportional to $2^{-258}\sqrt{|a-b|/L}$ yields exponential
rate at least $2^{-260}$ in that square root. The chosen
$\cs=2^{-1024}$ is smaller. This proves
\eqref{c:eq:cellgevrey} with the stated constants, including the
product inverse in~\eqref{c:eq:Rderivative}. The nonsingular contour
estimate similarly has a Euclidean rate at least $2^{-256}/L$,
again larger than the stated $\cs/L$.

For the nested-cutoff leakage, the high-momentum bound $2^{192}$
follows from $H\ge2^{99}I$ and $q_0=2^{-140}$, with a factor four
for the denominator remainder. The low-input constant after the
$h^2$ cancellation is at most $2^{232}J$. With
$K\le2L^{1/3}$, the Fourier-tail exponent satisfies
\[
 \frac1{512}\sqrt{q_0L/(4K)}
 \ge2^{-82}L^{1/3}.
\]
Together with the contraction $4^{-K}$ this proves the stated
leakage rate with $\cs=2^{-1024}$. The finite geometric sums of
contraction factors are below two. Hence no factor exponential
in $K$ is hidden in the leakage constant.

\subsection{Effective comparison and regularity constants}
The Fourier-alias error in~\eqref{c:eq:umklapperror}, divided by $h$,
is bounded by $2^{100}s^8e^{-s/512}$ with $s=h^{-1/2}$.
Its maximum occurs at $s=4096$ and is below $2^{196}$.
For~\eqref{c:eq:multipliererror}, the quartic contribution is at most
\[
 h^2 M_4|p|/2^{297/2}
 \le h\,2^{350-138-297/2}<2^{64}h;
\]
the cutoff tail is at most $2^{91}h$, and the eigenvector-multiplier
error from its Cauchy derivative is below $2^{70}h$.
Thus the displayed $2^{128}h$ multiplier bound and the
$2^{513}h$ operator bound leave room for the $2^{520}h$ comparison
used in the perturbation estimates. The scalar gap, reconstruction
and real normalization give the spatial-mode bound $2^{921}h$ in
\eqref{c:eq:modeapprox}. The bounded separated part contributes only
$2^{112}h^2$ before gap division.

For translations, the material derivative is at most $7\pi\varrho$,
and the local inverse norm is at most $4L^2$. The resolvent identity
therefore bounds its translation difference by $2^{20}L^4|a|$.
The bump factors add at most $2^{15}L|a|$, so $2^{32}L^4|a|$
bounds the translated compression. The spectral gap
$2^{-192}L^2$ then gives the normalized eigenvector difference
quotient below $2^{256}L^2$, and its full $H^1$ norm below
$2^{512}L^2$. This establishes the constants used for the cutoff
field's domain, without assuming ellipticity of curl--curl on the
whole vector space.

\subsection{Closing the fixed constant ledger}
All scalar, matrix, and derivative estimates just listed use at most
32 derivatives. Their product, inverse-derivative, Gaussian, and
finite-support summations are bounded by $M_{32}^{1024}$.
For an explicit check on the combinatorics, a derivative product is
summed with the binomial theorem; the resulting factors are bounded
by $2^j(j!)^4$, already present in $M_j$. The inverse recurrence
has at most $j$ terms at its $j$th step and is bounded by induction
by $M_j^4$. Cauchy division and the two cutoff products use at most
six such families; the pole/Gaussian remainder and the regular
integral use fewer than sixteen further products of their norms.
Even assigning 32 factors $M_{32}^{16}$ to each of these stages
stays below $M_{32}^{1024}$. The large support factors are the
explicit powers of $L$, or the finite counts $J,N_\Omega$, not
part of this fixed constant.

Finally
\[
 \log_2(M_{32}^{1024})
 =1024\bigl(4\log_2(32!)+16384\cdot33^2\bigr)
 <2^{35}<2^{40}=\log_2\Cs.
\]
Thus the common analytic coefficient $\Cs$ used in the interaction,
cell, leakage, and comparison estimates exceeds this fully specified
envelope. Whenever several such estimates are composed later, the
proof uses the displayed powers $\Cs^4,\Cs^5,\Cs^{32}$ rather
than silently replacing a product by $\Cs$.
Also $L\ge\Ls\ge\Cs^{64}$ makes every fixed smallness requirement
used above strict, for example $2^{520}h<2^{-200}$ and the errors
in the positive source numerator. The value $d\ge\Cs^8$ absorbs
the fixed large-distance thresholds in the Gaussian asymptotic.
No constant in this ledger is defined by the outcome of a
nonterminating search or an inaccessible exact eigenfunction.
